%% =====================================================================
%% chap_model.tex -- Master's thesis chapter:
%%   Kinematics, dynamics, and the parameter structure of the
%%   payload-adaptive model.
%%
%% FORMAT: single-column thesis chapter (\chapter / \section /
%% \subsection). Contains no \documentclass or preamble, so it drops
%% into any thesis template with %% =====================================================================
%% chap_model.tex -- Master's thesis chapter:
%%   Kinematics, dynamics, and the parameter structure of the
%%   payload-adaptive model.
%%
%% FORMAT: single-column thesis chapter (\chapter / \section /
%% \subsection). Contains no \documentclass or preamble, so it drops
%% into any thesis template with %% =====================================================================
%% chap_model.tex -- Master's thesis chapter:
%%   Kinematics, dynamics, and the parameter structure of the
%%   payload-adaptive model.
%%
%% FORMAT: single-column thesis chapter (\chapter / \section /
%% \subsection). Contains no \documentclass or preamble, so it drops
%% into any thesis template with %% =====================================================================
%% chap_model.tex -- Master's thesis chapter:
%%   Kinematics, dynamics, and the parameter structure of the
%%   payload-adaptive model.
%%
%% FORMAT: single-column thesis chapter (\chapter / \section /
%% \subsection). Contains no \documentclass or preamble, so it drops
%% into any thesis template with \input{chap_model}.
%%
%% REQUIRES in the preamble:
%%   amsmath, amssymb, booktabs, siunitx, bm
%%   (optional) \numberwithin{equation}{chapter}  for 3.1, 3.2, ...
%%
%% NOTE: the conference version of this material is sec_model.tex
%% (two-column IEEEtran, condensed). Keep both; they diverge on
%% purpose -- the thesis version carries the full derivations.
%% =====================================================================

% ---- Notation macros (move to the preamble if preferred) ------------
\providecommand{\dJ}{\Delta J}
\providecommand{\Jt}{J_t}
\providecommand{\Jb}{J_b}
\providecommand{\cvec}{\bm{c}}
\providecommand{\rp}{\bm{r}_p}
\providecommand{\rc}{\bm{r}_c}
\providecommand{\mpay}{\hat{m}_p}
\providecommand{\mb}{m_b}
\providecommand{\mt}{m_t}
\providecommand{\Tphys}{T_{\mathrm{phys}}}
\providecommand{\tauphys}{\bm{\tau}_{\mathrm{phys}}}
\providecommand{\tcom}{\bm{\tau}_{\mathrm{com}}}
\providecommand{\ev}{\bm{e}_3}

\chapter{Modeling and Parameter Structure}\label{ch:model}

A quadrotor that picks up a parcel does not merely become heavier. It
also becomes harder to rotate, and its centre of mass no longer lies on
the thrust axis. These are three distinct changes to the equations of
motion, and a controller that compensates only the first will fail in
ways that mass adaptation cannot explain.

This chapter derives the model that the estimator and the controller
share, and then makes explicit which of its parameters are identified
online and which are not. The central structural claim --- and the reason
the method needs so little sensing --- is that the three payload-induced
quantities are not three independent unknowns. Only the mass is a
decision variable of the estimator; the inertia increment and the
centre-of-mass offset are generated from it and from known contact
geometry by rigid-body algebra, at the cost of no additional degrees of
freedom.

Section~\ref{sec:model:frames} fixes frames and conventions,
Sections~\ref{sec:model:kin}--\ref{sec:model:dyn} derive the kinematics
and dynamics, and Section~\ref{sec:model:payload} derives the payload
parameterization. Sections~\ref{sec:model:phys}--\ref{sec:model:ctrl}
describe how the parameters are measured and consumed, and
Sections~\ref{sec:model:whynot}--\ref{sec:model:approx} defend the design
and declare its approximations.

\section{Frames, Notation, and Conventions}\label{sec:model:frames}

Two reference frames are used throughout. The world frame $\mathcal{W}$
is East--North--Up (ENU), so that gravity acts along $-z$. The body frame
$\mathcal{B}$ is Front--Left--Up (FLU) with its origin at the geometric
centre of the rotor disc --- \emph{not} at the centre of mass. This
choice is deliberate and consequential: the thrust vector passes through
the body origin by construction, so any offset between that origin and
the true centre of mass appears explicitly as a moment
(Section~\ref{sec:model:dyn}) rather than being absorbed silently into
the frame definition.

The state vector and the control input are
\begin{equation}
x = \bigl[\bm{p};\, \bm{v};\, \bm{q};\, \bm{\omega}\bigr] \in \mathbb{R}^{13},
\qquad
u = \bigl[T;\, \bm{\tau}\bigr] \in \mathbb{R}^{4},
\label{eq:state}
\end{equation}
where $\bm{p}$ and $\bm{v}$ are the position and velocity in
$\mathcal{W}$, $\bm{\omega}$ is the body rate in $\mathcal{B}$, $T$ is the
collective thrust along the body $z$ axis, and $\bm{\tau}$ is the body
torque. We abbreviate $\ev = [0,0,1]^{\!\top}$.

\subsection{Attitude representation}\label{sec:model:quat}

Attitude is represented by a unit quaternion under the \emph{Hamilton}
convention with the scalar part first,
\begin{equation}
\bm{q} = [q_w,\, q_x,\, q_y,\, q_z]^{\!\top},
\qquad \|\bm{q}\| = 1,
\label{eq:quatdef}
\end{equation}
encoding the rotation from $\mathcal{B}$ to $\mathcal{W}$. The
corresponding rotation matrix is
\begin{equation}
R(\bm{q}) =
\begin{bmatrix}
q_w^2 + q_x^2 - q_y^2 - q_z^2 & 2(q_xq_y - q_wq_z) & 2(q_xq_z + q_wq_y) \\[2pt]
2(q_xq_y + q_wq_z) & q_w^2 - q_x^2 + q_y^2 - q_z^2 & 2(q_yq_z - q_wq_x) \\[2pt]
2(q_xq_z - q_wq_y) & 2(q_yq_z + q_wq_x) & q_w^2 - q_x^2 - q_y^2 + q_z^2
\end{bmatrix}.
\label{eq:Rq}
\end{equation}

Three conventions have just been fixed, and all three matter. Hamilton
versus JPL determines the sign of the vector part in the quaternion
product; scalar-first versus scalar-last determines the layout of
\eqref{eq:Rq}; and resolving $\bm{\omega}$ in $\mathcal{B}$ rather than
$\mathcal{W}$ determines on which \emph{side} the product appears in the
kinematics \eqref{eq:quat} below. A mismatch in any one of them
integrates the attitude in the wrong direction while remaining
dimensionally consistent, which makes such errors unusually difficult to
detect. We state them here so that every later equation is unambiguous.

\section{Rigid-Body Kinematics}\label{sec:model:kin}

Translational kinematics is the definition of velocity,
\begin{equation}
\dot{\bm{p}} = \bm{v}.
\label{eq:kin}
\end{equation}

Attitude kinematics follows from the quaternion product of the attitude
with the pure quaternion formed from the body rate. Writing
$\bm{q} = (q_w, \bm{q}_v)$ with $\bm{q}_v = [q_x, q_y, q_z]^{\!\top}$,
and using the Hamilton product
$(a_w, \bm{a}_v) \otimes (b_w, \bm{b}_v) =
(a_w b_w - \bm{a}_v \!\cdot\! \bm{b}_v,\;
a_w \bm{b}_v + b_w \bm{a}_v + \bm{a}_v \!\times\! \bm{b}_v)$,
we obtain
\begin{equation}
\dot{\bm{q}} = \tfrac{1}{2}\, \bm{q} \otimes [0,\, \bm{\omega}]
= \tfrac{1}{2}
\begin{bmatrix}
-\bm{q}_v \cdot \bm{\omega} \\[2pt]
q_w \bm{\omega} + \bm{q}_v \times \bm{\omega}
\end{bmatrix}
= \tfrac{1}{2}\, \Xi(\bm{q})\, \bm{\omega},
\label{eq:quat}
\end{equation}
where the second equality expands the product and the third collects the
result into the $4 \times 3$ matrix
\begin{equation}
\Xi(\bm{q}) =
\begin{bmatrix}
-q_x & -q_y & -q_z \\
 q_w & -q_z &  q_y \\
 q_z &  q_w & -q_x \\
-q_y &  q_x &  q_w
\end{bmatrix}.
\label{eq:Xi}
\end{equation}
The two forms are algebraically identical; the implementation uses
\eqref{eq:Xi} because it is a plain matrix--vector product. Note that the
body rate enters on the \emph{right} of the product in \eqref{eq:quat}.
Had $\bm{\omega}$ been resolved in $\mathcal{W}$, the correct expression
would be $\dot{\bm{q}} = \tfrac{1}{2} [0, \bm{\omega}] \otimes \bm{q}$
instead, differing by a conjugation by $\bm{q}$.

Neither \eqref{eq:kin} nor \eqref{eq:quat} contains a mass, an inertia, or
any other physical parameter: kinematics is pure geometry. \emph{The
payload therefore enters the model only through the dynamics}, which
confines the entire parameter discussion of this chapter to the three
equations of Section~\ref{sec:model:dyn}.

\section{Rigid-Body Dynamics}\label{sec:model:dyn}

\subsection{Translational dynamics}\label{sec:model:trans}

Newton's second law in $\mathcal{W}$, with thrust rotated from the body
frame and a linear aerodynamic drag term, gives
\begin{equation}
\dot{\bm{v}} = \frac{1}{m}\Bigl( R(\bm{q})\, \ev\, T - k_d\, \bm{v} \Bigr) - g\, \ev,
\label{eq:trans}
\end{equation}
where $m$ is the \emph{composite} mass of airframe plus payload, $k_d$ is
a fixed drag coefficient, and $g$ is the gravitational acceleration.

Equation \eqref{eq:trans} is where the mass --- and \emph{only} the mass
--- enters, and it enters in exactly one role: as the scalar gain $1/m$
from thrust to acceleration. This single fact underwrites the
observability argument of the whole method. If a thrust value can be
obtained independently of the model, then the measured acceleration
determines $m$ uniquely by \eqref{eq:trans}, and it does so in ordinary
near-hover flight, with no dedicated excitation manoeuvre. The
qualification ``independently of the model'' is essential and is the
subject of Section~\ref{sec:model:phys}.

\subsection{Rotational dynamics}\label{sec:model:rot}

Let $\Jb = \mathrm{diag}(J_{xx}, J_{yy}, J_{zz})$ denote the
bare-airframe inertia, and let the payload raise the roll and pitch axes
by an increment $\dJ$. We write the resulting \emph{composite inertia} as
\begin{equation}
\Jt = \mathrm{diag}\bigl(J_{xx} + \dJ,\; J_{yy} + \dJ,\; J_{zz}\bigr),
\label{eq:Jt}
\end{equation}
the subscript matching the composite mass $\mt$ introduced in
Section~\ref{sec:model:payload}. Euler's equation for a rigid body,
written about the composite centre of mass and in the body frame, is then
\begin{equation}
\dot{\bm{\omega}} = \Jt^{-1}\Bigl( \bm{\tau} + \tcom(\cvec) - \bm{\omega} \times \Jt \bm{\omega} \Bigr).
\label{eq:rot}
\end{equation}
The term $\bm{\omega} \times \Jt \bm{\omega}$ is the usual gyroscopic
coupling. The term $\tcom$ is derived next.

That \eqref{eq:rot} is written \emph{about the composite centre of mass}
is not a stylistic choice: Euler's equation takes this form only when the
inertia tensor and the applied moments refer to the same point, and that
point is the centre of mass. Both $\Jt$ in Section~\ref{sec:model:payload}
and $\tcom$ below must therefore be computed about the composite CoM, not
about the body origin. Mixing the two reference points is a common and
silent modelling error.

\subsection{The thrust eccentricity moment}\label{sec:model:tcom}

Thrust acts at the rotor-disc centre, which by the frame convention of
Section~\ref{sec:model:frames} is the body origin. When a payload is
attached the composite centre of mass moves to some
$\rc = [c_x, c_y, c_z]^{\!\top} \neq \bm{0}$, and the thrust therefore no
longer passes through it. The moment about the CoM is the cross product
of the position of the application point \emph{relative to the CoM},
namely $-\rc$, with the force $\ev T$:
\begin{align}
\tcom(\cvec) &= (-\rc) \times \ev T
= -\begin{bmatrix} c_x \\ c_y \\ c_z \end{bmatrix}
\times \begin{bmatrix} 0 \\ 0 \\ T \end{bmatrix} \nonumber \\[4pt]
&= -\begin{bmatrix} c_y T - 0 \\ 0 - c_x T \\ 0 - 0 \end{bmatrix}
= \begin{bmatrix} -c_y T \\ +c_x T \\ 0 \end{bmatrix}.
\label{eq:tcom}
\end{align}

Two properties of \eqref{eq:tcom} deserve emphasis.

First, \textbf{$c_z$ cannot appear.} The third component of $\rc$
multiplies only terms that vanish, because a vertical CoM offset is
collinear with the body-$z$ thrust and the cross product of collinear
vectors is identically zero. Gravity, likewise, acts \emph{at} the centre
of mass and exerts no moment about it. The centre-of-mass offset is
therefore a two-parameter quantity \emph{by construction, not by
truncation} --- a distinction worth stating explicitly, since ``the
offset is three-dimensional, yet only two components are modelled'' is an
entirely natural objection to raise.

Second, \textbf{$\tcom$ is a constant moment in level hover.} It does not
average to zero, and it is not small: for the payload of
Section~\ref{sec:model:payload} it consumes roughly two thirds of the
available roll authority. Hovering on the spot with an off-centre load
therefore \emph{requires} a persistent trim torque. A model that omits
\eqref{eq:tcom} does not merely lose accuracy; it presents the controller
with a standing disturbance that has no explanation, which the controller
will then try to eliminate by manoeuvring. Section~\ref{sec:model:ctrl}
shows that the cost function must admit the same fact independently.

\subsection{Structure of the parameter entry points}\label{sec:model:entry}

Collecting \eqref{eq:trans}, \eqref{eq:rot} and \eqref{eq:tcom}, the model
has the property on which the rest of this thesis is built:

\begin{quote}
\emph{Each payload-dependent parameter has exactly one point of entry.}
The mass appears only as $1/m$ in the translational dynamics; the inertia
increment appears only inside $\Jt$ in the rotational dynamics; the
centre-of-mass offset appears only in the eccentricity moment.
\end{quote}

This separation is what makes the parameters individually attributable to
distinct physical channels --- the mass to the thrust--acceleration
relation, the inertia to angular acceleration, the offset to the
steady-state torque balance --- and it is exploited directly in
Sections~\ref{sec:model:mhe} and \ref{sec:model:whynot}.

\section{Payload-Induced Parameters}\label{sec:model:payload}

We now derive $\dJ$ and $\cvec$. The payload is modelled as a point mass
$\mpay$ rigidly attached at a known body-frame offset
$\rp = [r_x, r_y, r_z]^{\!\top}$ with $r_z < 0$ (the load hangs below the
airframe). Define the composite mass and the \emph{reduced mass}
\begin{equation}
\mt = \mb + \mpay,
\qquad
\mu = \frac{\mb\, \mpay}{\mb + \mpay} = \frac{\mb\, \mpay}{\mt},
\label{eq:mu}
\end{equation}
where $\mb$ is the bare-airframe mass.

\subsection{Centre-of-mass offset}\label{sec:model:com}

The composite centre of mass is the mass-weighted mean of the two
constituent positions. Taking the airframe centre of mass at the body
origin and the payload at $\rp$,
\begin{equation}
\rc = \frac{\mb \cdot \bm{0} + \mpay \cdot \rp}{\mt}
= \frac{\mpay}{\mt}\, \rp ,
\label{eq:rc}
\end{equation}
whose horizontal components are the pair required by \eqref{eq:tcom}:
\begin{equation}
\cvec = [c_x,\, c_y] = \frac{\mpay}{\mt}\, [r_x,\, r_y].
\label{eq:cxy}
\end{equation}

\subsection{Inertia increment and the reduced mass}\label{sec:model:paxis}

The inertia increment requires more care than the parallel-axis theorem
applied naively, because \emph{the centre of mass itself moves when the
payload is attached}. The theorem must be applied about the \emph{new}
centre of mass \eqref{eq:rc}, and both bodies contribute, since the
airframe no longer sits at the CoM either.

Consider a single axis and let $d$ denote the corresponding lever arm.
From \eqref{eq:rc}, the distances of the two bodies from the composite
centre of mass are
\begin{equation}
\underbrace{\bigl\| \bm{0} - \rc \bigr\| = \frac{\mpay}{\mt}\, d}_{\text{airframe}},
\qquad
\underbrace{\bigl\| \rp - \rc \bigr\| = \Bigl( 1 - \frac{\mpay}{\mt} \Bigr) d = \frac{\mb}{\mt}\, d}_{\text{payload}} .
\label{eq:dists}
\end{equation}
Note that the payload's distance is \emph{smaller} than $d$. Adding the
two parallel-axis contributions and simplifying,
\begin{align}
\Delta J^{(d)}
&= \mb \Bigl( \frac{\mpay}{\mt} d \Bigr)^{\!2}
 + \mpay \Bigl( \frac{\mb}{\mt} d \Bigr)^{\!2} \nonumber \\[4pt]
&= \frac{d^2}{\mt^2}\Bigl( \mb\, \mpay^2 + \mpay\, \mb^2 \Bigr)
 = \frac{d^2}{\mt^2}\, \mb\, \mpay \bigl( \mpay + \mb \bigr) \nonumber \\[4pt]
&= \frac{\mb\, \mpay}{\mt}\, d^2
 = \mu\, d^2 .
\label{eq:reduced}
\end{align}
The factor $(\mb + \mpay)$ produced by the sum cancels one power of $\mt$
in the denominator, and the two contributions collapse onto the reduced
mass \eqref{eq:mu} --- the same quantity that governs the two-body
problem in classical mechanics, and for the same reason: the inertia of a
rigid pair depends only on their \emph{relative} configuration.

Equation \eqref{eq:reduced} also identifies a natural error. Writing
$\Delta J^{(d)} = \mpay d^2$ --- treating the payload as orbiting a fixed
body origin, and forgetting both the airframe's own contribution and the
displacement of the CoM --- overestimates the increment by the factor
$\mt / \mb$, which is \SI{14.5}{\percent} at the payload considered
below. Two limiting cases confirm \eqref{eq:reduced}: as
$\mpay \!\ll\! \mb$ we recover $\mu \to \mpay$ (a light load barely moves
the CoM, and the naive expression becomes correct), while as
$\mpay \!\gg\! \mb$ we obtain $\mu \to \mb$ (a heavy load anchors the CoM
and it is the airframe that revolves about it).

Applying \eqref{eq:reduced} per axis with the appropriate arms gives
\begin{equation}
\dJ_{xx} = \mu\bigl(r_y^2 + r_z^2\bigr), \quad
\dJ_{yy} = \mu\bigl(r_x^2 + r_z^2\bigr), \quad
\dJ_{zz} = \mu\bigl(r_x^2 + r_y^2\bigr),
\label{eq:djaxes}
\end{equation}
and the model carries the roll/pitch mean as a single scalar,
\begin{equation}
\dJ = \tfrac{1}{2}\bigl(\dJ_{xx} + \dJ_{yy}\bigr)
    = \mu\Bigl( r_z^2 + \tfrac{1}{2}\bigl(r_x^2 + r_y^2\bigr) \Bigr).
\label{eq:dj}
\end{equation}
The consequences of using one scalar in place of the pair, and of
omitting $\dJ_{zz}$ from \eqref{eq:Jt}, are stated in
Section~\ref{sec:model:approx}.

\subsection{One unknown, three numbers}\label{sec:model:onescalar}

Equations \eqref{eq:cxy} and \eqref{eq:dj} produce all three
payload-dependent quantities --- $\dJ$, $c_x$, $c_y$ --- from a
\emph{single} unknown scalar $\mpay$, together with geometry $\rp$ that
is known from the gripper contact configuration. The rigid-attachment
constraint thus supplies two degrees of freedom that a naive augmentation
would otherwise have to identify from data. The estimator is not asked to
rediscover the parallel-axis theorem from measurement noise; it is given
the theorem and asked for one number.

This is the structural claim of the chapter, and
Table~\ref{tab:params} records the resulting parameter identities.
Section~\ref{sec:model:whynot} defends the choice against the obvious
alternative.

\begin{table}[t]
\caption{Model parameters, classified by identity. Only the mass is a
decision variable of the estimator; the inertia increment and the CoM
offset enter the estimator as \emph{known} runtime parameters, on the same
footing as the measured thrust and torque.}
\label{tab:params}
\centering
\footnotesize
\begin{tabular}{@{}l p{0.23\textwidth} c p{0.235\textwidth} p{0.17\textwidth}@{}}
\toprule
Symbol & Meaning & Dim. & Identity & Enters through \\
\midrule
$m$ & composite mass & 1 & \textbf{MHE decision variable} & \eqref{eq:trans}, \eqref{eq:hover} \\
$\dJ$ & roll/pitch inertia increment & 1 & derived, \eqref{eq:dj} & \eqref{eq:Jt}, \eqref{eq:gain} \\
$\cvec$ & CoM horizontal offset & 2 & derived \eqref{eq:cxy}, or inverted online \eqref{eq:conline} & \eqref{eq:tcom}, \eqref{eq:hover} \\
\midrule
$\mpay$ & payload mass (\emph{sole unknown}) & 1 & weak operational prior & \eqref{eq:cxy}, \eqref{eq:dj} \\
$\rp$ & payload offset in $\mathcal{B}$ & 3 & known (contact geometry) & \eqref{eq:cxy}, \eqref{eq:dj} \\
\midrule
$\mb$ & bare-airframe mass & 1 & calibration & \SI{2.0643}{\kg} \\
$\Jb$ & bare-airframe inertia & 3 & calibration & $\mathrm{diag}(0.0142,$ $0.0142, 0.0210)$ \\
$k_d$, $g$ & drag coefficient, gravity & 2 & constants & $0.05$, $9.81$ \\
\bottomrule
\end{tabular}
\end{table}

\subsection{Magnitudes on the experimental platform}\label{sec:model:magnitudes}

To fix ideas, consider a \SI{0.3}{\kg} payload attached at
$\rp = [0,\, 0.109,\, -0.47]\,$\si{\meter} on the platform of this work
($\mb = \SI{2.0643}{\kg}$, $J_{xx} = \SI{0.0142}{\kg\meter\squared}$,
roll/pitch torque authority \SI{0.5}{\newton\meter}). Equations
\eqref{eq:mu}, \eqref{eq:cxy} and \eqref{eq:dj} give
\begin{align}
\mu &= \SI{0.262}{\kg}, &
\dJ &= \SI{0.0594}{\kg\meter\squared}, \nonumber \\
\frac{J_{xx} + \dJ}{J_{xx}} &= 5.18, &
c_y &= \SI{13.8}{\milli\meter},
\label{eq:numbers}
\end{align}
and hence a constant hover trim torque
\begin{equation}
|c_y|\, \mt\, g = \SI{0.32}{\newton\meter} \approx \SI{64}{\percent}\ \text{of the roll authority.}
\label{eq:trimtorque}
\end{equation}

These numbers state the problem this thesis addresses more sharply than
any qualitative argument. A \SI{0.3}{\kg} parcel on a \SI{2.06}{\kg}
airframe is a \SI{15}{\percent} change of mass --- but simultaneously a
\emph{five-fold} change of roll inertia and a standing demand on two
thirds of the torque margin. Adapting the mass alone accounts for neither
of the latter two. Note also that the $r_z^2$ term contributes
\SI{97.4}{\percent} of $\dJ$ in \eqref{eq:dj}: the increment is dominated
by how far \emph{below} the airframe the load hangs, not by how far it is
off-centre.

\section{Model-Independent Reconstruction of the Physical Input}\label{sec:model:phys}

The thrust and torque appearing in \eqref{eq:trans} and \eqref{eq:rot}
must be known to the estimator. They are reconstructed from measured
rotor speeds $\omega_i$ and the known rotor positions $(x_i, y_i)$ via the
quadratic rotor model $F_i = k_f \omega_i^2$:
\begin{equation}
\Tphys = \sum_i F_i, \qquad
\tau_{\mathrm{roll}} = \sum_i y_i F_i, \qquad
\tau_{\mathrm{pitch}} = -\sum_i x_i F_i .
\label{eq:phys}
\end{equation}

It is worth being precise about why the \emph{commanded} thrust cannot be
used instead, since it is available for free and appears equivalent. Let
the controller compute its thrust command from an assumed mass
$\tilde{m}$, so that in hover $T_{\mathrm{cmd}} = \tilde{m} g$. Feeding
$T_{\mathrm{cmd}}$ into \eqref{eq:trans} and solving for the mass yields
$m = T_{\mathrm{cmd}} / g = \tilde{m}$ --- for \emph{any} $\tilde{m}$
whatsoever. The relation is satisfied identically, the residual carries no
information, and the mass is unobservable. The estimate would simply
reproduce whatever the controller already believed.

Equation \eqref{eq:phys} breaks this circularity because rotor speeds are
a physical measurement into which no model state enters. This is a
recurring principle in what follows: an adaptive loop must be fed by a
signal that is not itself a product of the adaptation.

\section{Estimation: The Mass as the Single Online Degree of Freedom}\label{sec:model:mhe}

\subsection{Augmented state and estimator formulation}\label{sec:model:mheform}

The moving-horizon estimator augments the mass into the state,
\begin{equation}
x_{\mathrm{aug}} = [x;\, m] \in \mathbb{R}^{14},
\qquad \dot m = 0,
\label{eq:xaug}
\end{equation}
holding it constant across the horizon so that a window of measurements
constrains a single value rather than a free trajectory. Process noise
$\bm{w} \in \mathbb{R}^{13}$ is admitted on the physical states only;
adding noise to the mass channel would allow it to drift within the
window and would dissolve precisely the constraint that makes the
estimate meaningful.

Over a horizon of $N$ stages the estimator solves
\begin{align}
\min_{x_0,\, \bm{w}} \quad
& \sum_{k=0}^{N-1} \Bigl( \bigl\| \bm{y}_k - h(x_k) \bigr\|^2_{R}
+ \bigl\| \bm{w}_k \bigr\|^2_{Q} \Bigr)
+ \bigl\| x_0 - \bar{x}_0 \bigr\|^2_{Q_0} \nonumber \\
\mathrm{s.t.} \quad
& x_{k+1} = f\bigl( x_k,\, \bm{w}_k;\, p_{\mathrm{MHE}} \bigr), \nonumber \\
& \underline{m} \le m \le \overline{m},
\label{eq:mhe}
\end{align}
where $f$ is the discretisation of \eqref{eq:kin}--\eqref{eq:rot},
$h$ selects the measured states, and the arrival cost
$\| x_0 - \bar{x}_0 \|^2_{Q_0}$ anchors the start of the window to the
estimate propagated from the previous solve.

The payload geometry does \emph{not} appear in the decision vector. It is
supplied as a runtime parameter, alongside the measured input:
\begin{equation}
p_{\mathrm{MHE}} = \bigl[\, T,\; \bm{\tau},\; \dJ,\; c_x,\; c_y \,\bigr] \in \mathbb{R}^{7}.
\label{eq:pmhe}
\end{equation}
Within the estimator, therefore, $\dJ$ and $\cvec$ have exactly the same
status as the measured thrust and torque: known inputs, not unknowns.

We use $N = 20$ stages at $dt = \SI{0.1}{\second}$, a \SI{2.0}{\second}
horizon, and bound the mass to
$[\SI{1.961}{\kg},\, \SI{5.0}{\kg}]$. The lower bound is the physical
floor $0.95\,\mb$: the vehicle cannot weigh less than its own airframe.
That this bound is a \emph{hard constraint on a decision variable} is a
concrete advantage of the moving-horizon formulation over a recursive
filter such as an extended Kalman filter, which can be biased away from
non-physical estimates but cannot be forbidden from producing them.

\subsection{Ground-truth-free inversion of the CoM offset}\label{sec:model:conline}

Where rotor speeds are available, the centre-of-mass offset need not come
from the prior at all. In steady hover the eccentricity moment
\eqref{eq:tcom} is balanced by the inner loop, so
$\dot{\bm{\omega}} \approx 0$ and $\bm{\omega} \approx \bm{0}$ reduce
\eqref{eq:rot} to $\bm{\tau} = -\tcom$. Substituting \eqref{eq:tcom} and
solving componentwise,
\begin{equation}
c_x = -\,\frac{\tau_{\mathrm{pitch}}^{\mathrm{phys}}}{\Tphys},
\qquad
c_y = +\,\frac{\tau_{\mathrm{roll}}^{\mathrm{phys}}}{\Tphys}.
\label{eq:conline}
\end{equation}

This is an algebraic inversion --- one division --- not an optimisation,
and it introduces no decision variable. Three properties make it
admissible within an adaptive loop.

\begin{enumerate}
\item \textbf{It reads no model state.} Both numerator and denominator
come from \eqref{eq:phys}. The inversion therefore cannot close an
estimate-to-model feedback path, which is the failure mode analysed in
Section~\ref{sec:model:whynot}.

\item \textbf{It is immune to thrust calibration.} Substituting
$F_i = k_f \omega_i^2$ into \eqref{eq:phys} and \eqref{eq:conline}, the
rotor coefficient cancels,
\begin{equation}
c_x = -\frac{-\sum_i x_i k_f \omega_i^2}{\sum_i k_f \omega_i^2}
    = \frac{\sum_i x_i \omega_i^2}{\sum_i \omega_i^2},
\label{eq:kfcancel}
\end{equation}
leaving a quantity that depends only on rotor geometry and the
\emph{relative} distribution of rotor speeds. Any method that consumes an
absolute thrust inherits the calibration error of $k_f$; this one does
not.

\item \textbf{Channels are assigned by identifiability.} The offset
$\cvec$ possesses an independent observable and is estimated online,
whereas $r_z$ is annihilated by the cross product in \eqref{eq:tcom} and
cannot be identified from the torque channel at all. It must remain a
prior. The split is dictated by the physics, not chosen for convenience.
\end{enumerate}

Two filtering stages precede use of \eqref{eq:conline}, and they address
different error mechanisms. A \emph{gate} rejects samples taken outside
steady flight ($\|\bm{\omega}\| < \SI{0.15}{\radian\per\second}$,
$\|\bm{v}_{xy}\| < \SI{0.20}{\meter\per\second}$), because during
manoeuvres the measured torque contains angular-acceleration terms and is
not a trim torque at all --- a \emph{systematic} error that no amount of
averaging would remove. An exponential moving average with time constant
\SI{2}{\second} then suppresses \emph{random} error, which is
substantial: the numerator of \eqref{eq:conline} is a differential-mode
combination of rotor forces, in which the common-mode signal cancels and
the noise does not. Low-pass filtering costs no information here because
$\cvec$ is quasi-static once the load is grasped --- its signal lies at
DC, while the rotor noise does not. The mass, by contrast, changes as a
step at grasp and release, and could not be treated this way; that is
precisely why $m$ is handled by the optimisation \eqref{eq:mhe} and
$\cvec$ by a first-order filter.

\section{Control: How the Parameters Re-centre the Cost}\label{sec:model:ctrl}

The nonlinear model predictive controller carries the parameters as
runtime values,
\begin{equation}
p_{\mathrm{NMPC}} = \bigl[\, x_{\mathrm{ref}}(13),\; m,\; \dJ,\; \cvec(2),\; \bm{d}(3) \,\bigr] \in \mathbb{R}^{20},
\label{eq:pnmpc}
\end{equation}
where $\bm{d}$ is a lumped translational disturbance used only by the
$\mathcal{L}_1$ comparison baseline and identically zero under the
proposed method. Over a prediction horizon of $N$ stages it solves
\begin{align}
\min_{u_{0:N-1}} \quad
& \sum_{k=0}^{N-1} \Bigl( \bigl\| e_{\mathrm{track}}(x_k, x_{\mathrm{ref},k}) \bigr\|^2_{Q}
+ \bigl\| u_k - u_{\mathrm{hover}} \bigr\|^2_{R} \Bigr) \nonumber \\
& \qquad + \bigl\| e_{\mathrm{track}}(x_N, x_{\mathrm{ref},N}) \bigr\|^2_{P} \nonumber \\
\mathrm{s.t.} \quad
& x_{k+1} = f(x_k, u_k;\, p_{\mathrm{NMPC}}), \nonumber \\
& \underline{T} \le T_k \le \overline{T}, \quad
  |\tau_{k,i}| \le \overline{\tau}_i ,
\label{eq:nmpc}
\end{align}
applying only $u_0$ and re-solving at the next step.

\subsection{Why the input baseline must track the parameters}\label{sec:model:hover}

The second term of \eqref{eq:nmpc} penalises the deviation of the input
from a \emph{baseline} $u_{\mathrm{hover}}$, and that baseline is
recomputed from the live parameters at every solve rather than fixed when
the solver is generated:
\begin{equation}
T_h = m\,g,
\qquad
u_{\mathrm{hover}} = \bigl[\, T_h,\; c_y T_h,\; -c_x T_h,\; 0 \,\bigr]^{\!\top}.
\label{eq:hover}
\end{equation}

Consider first why a baseline is needed at all. Penalising $\|u\|^2_R$
outright would treat the thrust required merely to remain airborne as an
expense to be minimised, biasing the solution toward insufficient thrust.
The baseline recentres the penalty on the input that sustains
equilibrium, so that $R$ penalises only genuine control effort.

Consider next what happens when the baseline is \emph{stale}. Suppose
$u_{\mathrm{hover}}$ is computed from the unloaded mass while the vehicle
carries a parcel. The cost then pulls the thrust persistently toward a
value below what equilibrium requires, and the optimum settles where the
marginal tracking cost balances the marginal baseline-deviation cost ---
at a nonzero steady-state error. The behaviour has a counter-intuitive
signature that identifies it in practice: \emph{increasing $R$ makes the
error worse}, because $R$ is precisely the strength with which the
incorrect baseline pulls. No amount of weight tuning removes a bias of
this kind; only correcting the baseline does.

The torque entries of \eqref{eq:hover} are exactly $-\tcom$ from
\eqref{eq:tcom}, and they exist for the same reason. Hovering with an
eccentric load requires a constant trim torque
(Section~\ref{sec:model:tcom}); a cost that does not admit this pulls the
torque toward zero and produces the same class of steady-state bias in
the attitude channel. The two equations are complementary and neither
suffices alone: \eqref{eq:tcom} teaches the \emph{model} that the moment
exists, while \eqref{eq:hover} teaches the \emph{objective} that it is
necessary.

\subsection{Inner-loop bandwidth}\label{sec:model:gain}

Model consistency in the outer loop cannot restore inner-loop bandwidth.
The rate controller running on the flight controller is tuned for the
bare airframe; a payload that multiplies the roll and pitch inertia
lowers its effective bandwidth, and for heavy eccentric loads the
resulting attitude oscillation grows until the torque saturates. We
therefore rescale the rate gains at attachment by the inertia ratio,
\begin{equation}
K \leftarrow K_{\mathrm{nom}} \cdot
\min\!\Bigl( \frac{J_{xx} + \dJ}{J_{xx}},\; 5.0 \Bigr).
\label{eq:gain}
\end{equation}

The cap is not incidental. From \eqref{eq:numbers}, a \SI{0.3}{\kg}
payload already produces a ratio of $5.18$, so that payload sits at the
ceiling of what inner-loop retuning alone can absorb on this airframe.
Together with the torque budget \eqref{eq:trimtorque}, this defines a
feasibility boundary on payload mass and eccentricity that is
characterised experimentally in a later chapter.

\section{Why the Inertia and CoM Are Not Augmented into the State}\label{sec:model:whynot}

Augmenting $\dJ$ and $\cvec$ into the estimator's decision vector
\eqref{eq:xaug} --- yielding a seventeen-dimensional augmented state and
four online parameters instead of one --- is the obvious alternative, and
it is the design most readers will expect. We argue against it on four
grounds.

\paragraph{Asymmetric observability} The mass is excited directly by the
thrust--acceleration relation \eqref{eq:trans} and is observable in
ordinary near-hover flight, as argued in
Section~\ref{sec:model:trans}. The inertia increment acts only through
$\dot{\bm{\omega}}$ in \eqref{eq:rot}, and a delivery task --- spent
largely in near-hover, low-rate flight --- excites that channel weakly.
Admitting $\dJ$ as a free variable therefore permits the optimiser to
explain measurement noise using a nearly unobservable quantity, which
degrades rather than improves the estimate.

\paragraph{Competition for a single residual} Deriving the geometry from
the \emph{instantaneous} mass estimate closes a feedback path,
\begin{equation*}
m_{\mathrm{est}} \;\longrightarrow\; \dJ,\, \cvec
\;\longrightarrow\; \text{attitude model}
\;\longrightarrow\; m_{\mathrm{est}},
\end{equation*}
in which two quantities compete to explain the same residual. This
bootstrap was observed to fail structurally on light payloads: before
lift-off the mass estimate reads low \emph{legitimately}, since the load
still rests on the ground; that value is consumed as ``no geometry
present''; the geometry channel closes; and after lift-off the model is
mismatched and the estimate remains pinned at its lower bound. The remedy
was to sever the path --- letting $\mpay$ read a fixed operational prior
and never the running estimate --- not to add further free parameters.

\paragraph{Economy of parameterisation} Equations \eqref{eq:cxy} and
\eqref{eq:dj} obtain three numbers from one scalar. The two degrees of
freedom thereby saved are not discarded information: they are recovered
from a constraint, rigid attachment, that holds exactly.

\paragraph{The remaining prior may be weak} The single scalar $\mpay$
need not be accurate. With the geometry prior deliberately set
\SI{33}{\percent} wrong, the mass estimate still converges to within
\SI{0.3}{\percent}, so a nominal parcel mass --- a quantity a delivery
operator knows anyway --- suffices, and no ground-truth measurement of
the attachment is required. Where rotor speeds are available,
\eqref{eq:conline} removes the prior from the $\cvec$ path entirely,
leaving it to influence only $\dJ$.

\section{Declared Modelling Approximations}\label{sec:model:approx}

Three points are stated explicitly here rather than left to be
discovered.

\paragraph{Anisotropy of the inertia increment} The single scalar
\eqref{eq:dj} replaces the pair $(\dJ_{xx}, \dJ_{yy})$ of
\eqref{eq:djaxes} by their arithmetic mean. The resulting error is
$\tfrac{1}{2}\mu\,|r_x^2 - r_y^2|$, which vanishes exactly when
$r_x = r_y$ and remains below \SI{1}{\percent} of $J_{xx} + \dJ$ at the
offsets used in this work.

\paragraph{The yaw increment is neglected} The third expression in
\eqref{eq:djaxes}, $\dJ_{zz} = \mu(r_x^2 + r_y^2)$, is not added to
$J_{zz}$ in \eqref{eq:Jt}. At \SI{0.3}{\kg} and
$r_{xy} = \SI{0.109}{\meter}$ it amounts to
\SI{0.0031}{\kg\meter\squared}, or \SI{14.8}{\percent} of $J_{zz}$. This
is a known and non-negligible approximation, and it is retained only
because the yaw channel neither carries payload compensation nor
participates in the inversion \eqref{eq:conline}, so it does not affect
the results reported later. Removing it would require substituting
$J_{zz} + \mu(r_x^2 + r_y^2)$ into \eqref{eq:Jt} and would introduce no
new estimated quantity --- the correction is available at zero cost in
identification, and is omitted only for consistency with the
implementation as evaluated.

\paragraph{The vertical CoM offset is absent by construction} As shown in
Section~\ref{sec:model:tcom}, $c_z$ drops out of \eqref{eq:tcom} as an
exact algebraic consequence of collinearity, not as a truncation. It is
listed among the approximations only to forestall the opposite reading.
.
%%
%% REQUIRES in the preamble:
%%   amsmath, amssymb, booktabs, siunitx, bm
%%   (optional) \numberwithin{equation}{chapter}  for 3.1, 3.2, ...
%%
%% NOTE: the conference version of this material is sec_model.tex
%% (two-column IEEEtran, condensed). Keep both; they diverge on
%% purpose -- the thesis version carries the full derivations.
%% =====================================================================

% ---- Notation macros (move to the preamble if preferred) ------------
\providecommand{\dJ}{\Delta J}
\providecommand{\Jt}{J_t}
\providecommand{\Jb}{J_b}
\providecommand{\cvec}{\bm{c}}
\providecommand{\rp}{\bm{r}_p}
\providecommand{\rc}{\bm{r}_c}
\providecommand{\mpay}{\hat{m}_p}
\providecommand{\mb}{m_b}
\providecommand{\mt}{m_t}
\providecommand{\Tphys}{T_{\mathrm{phys}}}
\providecommand{\tauphys}{\bm{\tau}_{\mathrm{phys}}}
\providecommand{\tcom}{\bm{\tau}_{\mathrm{com}}}
\providecommand{\ev}{\bm{e}_3}

\chapter{Modeling and Parameter Structure}\label{ch:model}

A quadrotor that picks up a parcel does not merely become heavier. It
also becomes harder to rotate, and its centre of mass no longer lies on
the thrust axis. These are three distinct changes to the equations of
motion, and a controller that compensates only the first will fail in
ways that mass adaptation cannot explain.

This chapter derives the model that the estimator and the controller
share, and then makes explicit which of its parameters are identified
online and which are not. The central structural claim --- and the reason
the method needs so little sensing --- is that the three payload-induced
quantities are not three independent unknowns. Only the mass is a
decision variable of the estimator; the inertia increment and the
centre-of-mass offset are generated from it and from known contact
geometry by rigid-body algebra, at the cost of no additional degrees of
freedom.

Section~\ref{sec:model:frames} fixes frames and conventions,
Sections~\ref{sec:model:kin}--\ref{sec:model:dyn} derive the kinematics
and dynamics, and Section~\ref{sec:model:payload} derives the payload
parameterization. Sections~\ref{sec:model:phys}--\ref{sec:model:ctrl}
describe how the parameters are measured and consumed, and
Sections~\ref{sec:model:whynot}--\ref{sec:model:approx} defend the design
and declare its approximations.

\section{Frames, Notation, and Conventions}\label{sec:model:frames}

Two reference frames are used throughout. The world frame $\mathcal{W}$
is East--North--Up (ENU), so that gravity acts along $-z$. The body frame
$\mathcal{B}$ is Front--Left--Up (FLU) with its origin at the geometric
centre of the rotor disc --- \emph{not} at the centre of mass. This
choice is deliberate and consequential: the thrust vector passes through
the body origin by construction, so any offset between that origin and
the true centre of mass appears explicitly as a moment
(Section~\ref{sec:model:dyn}) rather than being absorbed silently into
the frame definition.

The state vector and the control input are
\begin{equation}
x = \bigl[\bm{p};\, \bm{v};\, \bm{q};\, \bm{\omega}\bigr] \in \mathbb{R}^{13},
\qquad
u = \bigl[T;\, \bm{\tau}\bigr] \in \mathbb{R}^{4},
\label{eq:state}
\end{equation}
where $\bm{p}$ and $\bm{v}$ are the position and velocity in
$\mathcal{W}$, $\bm{\omega}$ is the body rate in $\mathcal{B}$, $T$ is the
collective thrust along the body $z$ axis, and $\bm{\tau}$ is the body
torque. We abbreviate $\ev = [0,0,1]^{\!\top}$.

\subsection{Attitude representation}\label{sec:model:quat}

Attitude is represented by a unit quaternion under the \emph{Hamilton}
convention with the scalar part first,
\begin{equation}
\bm{q} = [q_w,\, q_x,\, q_y,\, q_z]^{\!\top},
\qquad \|\bm{q}\| = 1,
\label{eq:quatdef}
\end{equation}
encoding the rotation from $\mathcal{B}$ to $\mathcal{W}$. The
corresponding rotation matrix is
\begin{equation}
R(\bm{q}) =
\begin{bmatrix}
q_w^2 + q_x^2 - q_y^2 - q_z^2 & 2(q_xq_y - q_wq_z) & 2(q_xq_z + q_wq_y) \\[2pt]
2(q_xq_y + q_wq_z) & q_w^2 - q_x^2 + q_y^2 - q_z^2 & 2(q_yq_z - q_wq_x) \\[2pt]
2(q_xq_z - q_wq_y) & 2(q_yq_z + q_wq_x) & q_w^2 - q_x^2 - q_y^2 + q_z^2
\end{bmatrix}.
\label{eq:Rq}
\end{equation}

Three conventions have just been fixed, and all three matter. Hamilton
versus JPL determines the sign of the vector part in the quaternion
product; scalar-first versus scalar-last determines the layout of
\eqref{eq:Rq}; and resolving $\bm{\omega}$ in $\mathcal{B}$ rather than
$\mathcal{W}$ determines on which \emph{side} the product appears in the
kinematics \eqref{eq:quat} below. A mismatch in any one of them
integrates the attitude in the wrong direction while remaining
dimensionally consistent, which makes such errors unusually difficult to
detect. We state them here so that every later equation is unambiguous.

\section{Rigid-Body Kinematics}\label{sec:model:kin}

Translational kinematics is the definition of velocity,
\begin{equation}
\dot{\bm{p}} = \bm{v}.
\label{eq:kin}
\end{equation}

Attitude kinematics follows from the quaternion product of the attitude
with the pure quaternion formed from the body rate. Writing
$\bm{q} = (q_w, \bm{q}_v)$ with $\bm{q}_v = [q_x, q_y, q_z]^{\!\top}$,
and using the Hamilton product
$(a_w, \bm{a}_v) \otimes (b_w, \bm{b}_v) =
(a_w b_w - \bm{a}_v \!\cdot\! \bm{b}_v,\;
a_w \bm{b}_v + b_w \bm{a}_v + \bm{a}_v \!\times\! \bm{b}_v)$,
we obtain
\begin{equation}
\dot{\bm{q}} = \tfrac{1}{2}\, \bm{q} \otimes [0,\, \bm{\omega}]
= \tfrac{1}{2}
\begin{bmatrix}
-\bm{q}_v \cdot \bm{\omega} \\[2pt]
q_w \bm{\omega} + \bm{q}_v \times \bm{\omega}
\end{bmatrix}
= \tfrac{1}{2}\, \Xi(\bm{q})\, \bm{\omega},
\label{eq:quat}
\end{equation}
where the second equality expands the product and the third collects the
result into the $4 \times 3$ matrix
\begin{equation}
\Xi(\bm{q}) =
\begin{bmatrix}
-q_x & -q_y & -q_z \\
 q_w & -q_z &  q_y \\
 q_z &  q_w & -q_x \\
-q_y &  q_x &  q_w
\end{bmatrix}.
\label{eq:Xi}
\end{equation}
The two forms are algebraically identical; the implementation uses
\eqref{eq:Xi} because it is a plain matrix--vector product. Note that the
body rate enters on the \emph{right} of the product in \eqref{eq:quat}.
Had $\bm{\omega}$ been resolved in $\mathcal{W}$, the correct expression
would be $\dot{\bm{q}} = \tfrac{1}{2} [0, \bm{\omega}] \otimes \bm{q}$
instead, differing by a conjugation by $\bm{q}$.

Neither \eqref{eq:kin} nor \eqref{eq:quat} contains a mass, an inertia, or
any other physical parameter: kinematics is pure geometry. \emph{The
payload therefore enters the model only through the dynamics}, which
confines the entire parameter discussion of this chapter to the three
equations of Section~\ref{sec:model:dyn}.

\section{Rigid-Body Dynamics}\label{sec:model:dyn}

\subsection{Translational dynamics}\label{sec:model:trans}

Newton's second law in $\mathcal{W}$, with thrust rotated from the body
frame and a linear aerodynamic drag term, gives
\begin{equation}
\dot{\bm{v}} = \frac{1}{m}\Bigl( R(\bm{q})\, \ev\, T - k_d\, \bm{v} \Bigr) - g\, \ev,
\label{eq:trans}
\end{equation}
where $m$ is the \emph{composite} mass of airframe plus payload, $k_d$ is
a fixed drag coefficient, and $g$ is the gravitational acceleration.

Equation \eqref{eq:trans} is where the mass --- and \emph{only} the mass
--- enters, and it enters in exactly one role: as the scalar gain $1/m$
from thrust to acceleration. This single fact underwrites the
observability argument of the whole method. If a thrust value can be
obtained independently of the model, then the measured acceleration
determines $m$ uniquely by \eqref{eq:trans}, and it does so in ordinary
near-hover flight, with no dedicated excitation manoeuvre. The
qualification ``independently of the model'' is essential and is the
subject of Section~\ref{sec:model:phys}.

\subsection{Rotational dynamics}\label{sec:model:rot}

Let $\Jb = \mathrm{diag}(J_{xx}, J_{yy}, J_{zz})$ denote the
bare-airframe inertia, and let the payload raise the roll and pitch axes
by an increment $\dJ$. We write the resulting \emph{composite inertia} as
\begin{equation}
\Jt = \mathrm{diag}\bigl(J_{xx} + \dJ,\; J_{yy} + \dJ,\; J_{zz}\bigr),
\label{eq:Jt}
\end{equation}
the subscript matching the composite mass $\mt$ introduced in
Section~\ref{sec:model:payload}. Euler's equation for a rigid body,
written about the composite centre of mass and in the body frame, is then
\begin{equation}
\dot{\bm{\omega}} = \Jt^{-1}\Bigl( \bm{\tau} + \tcom(\cvec) - \bm{\omega} \times \Jt \bm{\omega} \Bigr).
\label{eq:rot}
\end{equation}
The term $\bm{\omega} \times \Jt \bm{\omega}$ is the usual gyroscopic
coupling. The term $\tcom$ is derived next.

That \eqref{eq:rot} is written \emph{about the composite centre of mass}
is not a stylistic choice: Euler's equation takes this form only when the
inertia tensor and the applied moments refer to the same point, and that
point is the centre of mass. Both $\Jt$ in Section~\ref{sec:model:payload}
and $\tcom$ below must therefore be computed about the composite CoM, not
about the body origin. Mixing the two reference points is a common and
silent modelling error.

\subsection{The thrust eccentricity moment}\label{sec:model:tcom}

Thrust acts at the rotor-disc centre, which by the frame convention of
Section~\ref{sec:model:frames} is the body origin. When a payload is
attached the composite centre of mass moves to some
$\rc = [c_x, c_y, c_z]^{\!\top} \neq \bm{0}$, and the thrust therefore no
longer passes through it. The moment about the CoM is the cross product
of the position of the application point \emph{relative to the CoM},
namely $-\rc$, with the force $\ev T$:
\begin{align}
\tcom(\cvec) &= (-\rc) \times \ev T
= -\begin{bmatrix} c_x \\ c_y \\ c_z \end{bmatrix}
\times \begin{bmatrix} 0 \\ 0 \\ T \end{bmatrix} \nonumber \\[4pt]
&= -\begin{bmatrix} c_y T - 0 \\ 0 - c_x T \\ 0 - 0 \end{bmatrix}
= \begin{bmatrix} -c_y T \\ +c_x T \\ 0 \end{bmatrix}.
\label{eq:tcom}
\end{align}

Two properties of \eqref{eq:tcom} deserve emphasis.

First, \textbf{$c_z$ cannot appear.} The third component of $\rc$
multiplies only terms that vanish, because a vertical CoM offset is
collinear with the body-$z$ thrust and the cross product of collinear
vectors is identically zero. Gravity, likewise, acts \emph{at} the centre
of mass and exerts no moment about it. The centre-of-mass offset is
therefore a two-parameter quantity \emph{by construction, not by
truncation} --- a distinction worth stating explicitly, since ``the
offset is three-dimensional, yet only two components are modelled'' is an
entirely natural objection to raise.

Second, \textbf{$\tcom$ is a constant moment in level hover.} It does not
average to zero, and it is not small: for the payload of
Section~\ref{sec:model:payload} it consumes roughly two thirds of the
available roll authority. Hovering on the spot with an off-centre load
therefore \emph{requires} a persistent trim torque. A model that omits
\eqref{eq:tcom} does not merely lose accuracy; it presents the controller
with a standing disturbance that has no explanation, which the controller
will then try to eliminate by manoeuvring. Section~\ref{sec:model:ctrl}
shows that the cost function must admit the same fact independently.

\subsection{Structure of the parameter entry points}\label{sec:model:entry}

Collecting \eqref{eq:trans}, \eqref{eq:rot} and \eqref{eq:tcom}, the model
has the property on which the rest of this thesis is built:

\begin{quote}
\emph{Each payload-dependent parameter has exactly one point of entry.}
The mass appears only as $1/m$ in the translational dynamics; the inertia
increment appears only inside $\Jt$ in the rotational dynamics; the
centre-of-mass offset appears only in the eccentricity moment.
\end{quote}

This separation is what makes the parameters individually attributable to
distinct physical channels --- the mass to the thrust--acceleration
relation, the inertia to angular acceleration, the offset to the
steady-state torque balance --- and it is exploited directly in
Sections~\ref{sec:model:mhe} and \ref{sec:model:whynot}.

\section{Payload-Induced Parameters}\label{sec:model:payload}

We now derive $\dJ$ and $\cvec$. The payload is modelled as a point mass
$\mpay$ rigidly attached at a known body-frame offset
$\rp = [r_x, r_y, r_z]^{\!\top}$ with $r_z < 0$ (the load hangs below the
airframe). Define the composite mass and the \emph{reduced mass}
\begin{equation}
\mt = \mb + \mpay,
\qquad
\mu = \frac{\mb\, \mpay}{\mb + \mpay} = \frac{\mb\, \mpay}{\mt},
\label{eq:mu}
\end{equation}
where $\mb$ is the bare-airframe mass.

\subsection{Centre-of-mass offset}\label{sec:model:com}

The composite centre of mass is the mass-weighted mean of the two
constituent positions. Taking the airframe centre of mass at the body
origin and the payload at $\rp$,
\begin{equation}
\rc = \frac{\mb \cdot \bm{0} + \mpay \cdot \rp}{\mt}
= \frac{\mpay}{\mt}\, \rp ,
\label{eq:rc}
\end{equation}
whose horizontal components are the pair required by \eqref{eq:tcom}:
\begin{equation}
\cvec = [c_x,\, c_y] = \frac{\mpay}{\mt}\, [r_x,\, r_y].
\label{eq:cxy}
\end{equation}

\subsection{Inertia increment and the reduced mass}\label{sec:model:paxis}

The inertia increment requires more care than the parallel-axis theorem
applied naively, because \emph{the centre of mass itself moves when the
payload is attached}. The theorem must be applied about the \emph{new}
centre of mass \eqref{eq:rc}, and both bodies contribute, since the
airframe no longer sits at the CoM either.

Consider a single axis and let $d$ denote the corresponding lever arm.
From \eqref{eq:rc}, the distances of the two bodies from the composite
centre of mass are
\begin{equation}
\underbrace{\bigl\| \bm{0} - \rc \bigr\| = \frac{\mpay}{\mt}\, d}_{\text{airframe}},
\qquad
\underbrace{\bigl\| \rp - \rc \bigr\| = \Bigl( 1 - \frac{\mpay}{\mt} \Bigr) d = \frac{\mb}{\mt}\, d}_{\text{payload}} .
\label{eq:dists}
\end{equation}
Note that the payload's distance is \emph{smaller} than $d$. Adding the
two parallel-axis contributions and simplifying,
\begin{align}
\Delta J^{(d)}
&= \mb \Bigl( \frac{\mpay}{\mt} d \Bigr)^{\!2}
 + \mpay \Bigl( \frac{\mb}{\mt} d \Bigr)^{\!2} \nonumber \\[4pt]
&= \frac{d^2}{\mt^2}\Bigl( \mb\, \mpay^2 + \mpay\, \mb^2 \Bigr)
 = \frac{d^2}{\mt^2}\, \mb\, \mpay \bigl( \mpay + \mb \bigr) \nonumber \\[4pt]
&= \frac{\mb\, \mpay}{\mt}\, d^2
 = \mu\, d^2 .
\label{eq:reduced}
\end{align}
The factor $(\mb + \mpay)$ produced by the sum cancels one power of $\mt$
in the denominator, and the two contributions collapse onto the reduced
mass \eqref{eq:mu} --- the same quantity that governs the two-body
problem in classical mechanics, and for the same reason: the inertia of a
rigid pair depends only on their \emph{relative} configuration.

Equation \eqref{eq:reduced} also identifies a natural error. Writing
$\Delta J^{(d)} = \mpay d^2$ --- treating the payload as orbiting a fixed
body origin, and forgetting both the airframe's own contribution and the
displacement of the CoM --- overestimates the increment by the factor
$\mt / \mb$, which is \SI{14.5}{\percent} at the payload considered
below. Two limiting cases confirm \eqref{eq:reduced}: as
$\mpay \!\ll\! \mb$ we recover $\mu \to \mpay$ (a light load barely moves
the CoM, and the naive expression becomes correct), while as
$\mpay \!\gg\! \mb$ we obtain $\mu \to \mb$ (a heavy load anchors the CoM
and it is the airframe that revolves about it).

Applying \eqref{eq:reduced} per axis with the appropriate arms gives
\begin{equation}
\dJ_{xx} = \mu\bigl(r_y^2 + r_z^2\bigr), \quad
\dJ_{yy} = \mu\bigl(r_x^2 + r_z^2\bigr), \quad
\dJ_{zz} = \mu\bigl(r_x^2 + r_y^2\bigr),
\label{eq:djaxes}
\end{equation}
and the model carries the roll/pitch mean as a single scalar,
\begin{equation}
\dJ = \tfrac{1}{2}\bigl(\dJ_{xx} + \dJ_{yy}\bigr)
    = \mu\Bigl( r_z^2 + \tfrac{1}{2}\bigl(r_x^2 + r_y^2\bigr) \Bigr).
\label{eq:dj}
\end{equation}
The consequences of using one scalar in place of the pair, and of
omitting $\dJ_{zz}$ from \eqref{eq:Jt}, are stated in
Section~\ref{sec:model:approx}.

\subsection{One unknown, three numbers}\label{sec:model:onescalar}

Equations \eqref{eq:cxy} and \eqref{eq:dj} produce all three
payload-dependent quantities --- $\dJ$, $c_x$, $c_y$ --- from a
\emph{single} unknown scalar $\mpay$, together with geometry $\rp$ that
is known from the gripper contact configuration. The rigid-attachment
constraint thus supplies two degrees of freedom that a naive augmentation
would otherwise have to identify from data. The estimator is not asked to
rediscover the parallel-axis theorem from measurement noise; it is given
the theorem and asked for one number.

This is the structural claim of the chapter, and
Table~\ref{tab:params} records the resulting parameter identities.
Section~\ref{sec:model:whynot} defends the choice against the obvious
alternative.

\begin{table}[t]
\caption{Model parameters, classified by identity. Only the mass is a
decision variable of the estimator; the inertia increment and the CoM
offset enter the estimator as \emph{known} runtime parameters, on the same
footing as the measured thrust and torque.}
\label{tab:params}
\centering
\footnotesize
\begin{tabular}{@{}l p{0.23\textwidth} c p{0.235\textwidth} p{0.17\textwidth}@{}}
\toprule
Symbol & Meaning & Dim. & Identity & Enters through \\
\midrule
$m$ & composite mass & 1 & \textbf{MHE decision variable} & \eqref{eq:trans}, \eqref{eq:hover} \\
$\dJ$ & roll/pitch inertia increment & 1 & derived, \eqref{eq:dj} & \eqref{eq:Jt}, \eqref{eq:gain} \\
$\cvec$ & CoM horizontal offset & 2 & derived \eqref{eq:cxy}, or inverted online \eqref{eq:conline} & \eqref{eq:tcom}, \eqref{eq:hover} \\
\midrule
$\mpay$ & payload mass (\emph{sole unknown}) & 1 & weak operational prior & \eqref{eq:cxy}, \eqref{eq:dj} \\
$\rp$ & payload offset in $\mathcal{B}$ & 3 & known (contact geometry) & \eqref{eq:cxy}, \eqref{eq:dj} \\
\midrule
$\mb$ & bare-airframe mass & 1 & calibration & \SI{2.0643}{\kg} \\
$\Jb$ & bare-airframe inertia & 3 & calibration & $\mathrm{diag}(0.0142,$ $0.0142, 0.0210)$ \\
$k_d$, $g$ & drag coefficient, gravity & 2 & constants & $0.05$, $9.81$ \\
\bottomrule
\end{tabular}
\end{table}

\subsection{Magnitudes on the experimental platform}\label{sec:model:magnitudes}

To fix ideas, consider a \SI{0.3}{\kg} payload attached at
$\rp = [0,\, 0.109,\, -0.47]\,$\si{\meter} on the platform of this work
($\mb = \SI{2.0643}{\kg}$, $J_{xx} = \SI{0.0142}{\kg\meter\squared}$,
roll/pitch torque authority \SI{0.5}{\newton\meter}). Equations
\eqref{eq:mu}, \eqref{eq:cxy} and \eqref{eq:dj} give
\begin{align}
\mu &= \SI{0.262}{\kg}, &
\dJ &= \SI{0.0594}{\kg\meter\squared}, \nonumber \\
\frac{J_{xx} + \dJ}{J_{xx}} &= 5.18, &
c_y &= \SI{13.8}{\milli\meter},
\label{eq:numbers}
\end{align}
and hence a constant hover trim torque
\begin{equation}
|c_y|\, \mt\, g = \SI{0.32}{\newton\meter} \approx \SI{64}{\percent}\ \text{of the roll authority.}
\label{eq:trimtorque}
\end{equation}

These numbers state the problem this thesis addresses more sharply than
any qualitative argument. A \SI{0.3}{\kg} parcel on a \SI{2.06}{\kg}
airframe is a \SI{15}{\percent} change of mass --- but simultaneously a
\emph{five-fold} change of roll inertia and a standing demand on two
thirds of the torque margin. Adapting the mass alone accounts for neither
of the latter two. Note also that the $r_z^2$ term contributes
\SI{97.4}{\percent} of $\dJ$ in \eqref{eq:dj}: the increment is dominated
by how far \emph{below} the airframe the load hangs, not by how far it is
off-centre.

\section{Model-Independent Reconstruction of the Physical Input}\label{sec:model:phys}

The thrust and torque appearing in \eqref{eq:trans} and \eqref{eq:rot}
must be known to the estimator. They are reconstructed from measured
rotor speeds $\omega_i$ and the known rotor positions $(x_i, y_i)$ via the
quadratic rotor model $F_i = k_f \omega_i^2$:
\begin{equation}
\Tphys = \sum_i F_i, \qquad
\tau_{\mathrm{roll}} = \sum_i y_i F_i, \qquad
\tau_{\mathrm{pitch}} = -\sum_i x_i F_i .
\label{eq:phys}
\end{equation}

It is worth being precise about why the \emph{commanded} thrust cannot be
used instead, since it is available for free and appears equivalent. Let
the controller compute its thrust command from an assumed mass
$\tilde{m}$, so that in hover $T_{\mathrm{cmd}} = \tilde{m} g$. Feeding
$T_{\mathrm{cmd}}$ into \eqref{eq:trans} and solving for the mass yields
$m = T_{\mathrm{cmd}} / g = \tilde{m}$ --- for \emph{any} $\tilde{m}$
whatsoever. The relation is satisfied identically, the residual carries no
information, and the mass is unobservable. The estimate would simply
reproduce whatever the controller already believed.

Equation \eqref{eq:phys} breaks this circularity because rotor speeds are
a physical measurement into which no model state enters. This is a
recurring principle in what follows: an adaptive loop must be fed by a
signal that is not itself a product of the adaptation.

\section{Estimation: The Mass as the Single Online Degree of Freedom}\label{sec:model:mhe}

\subsection{Augmented state and estimator formulation}\label{sec:model:mheform}

The moving-horizon estimator augments the mass into the state,
\begin{equation}
x_{\mathrm{aug}} = [x;\, m] \in \mathbb{R}^{14},
\qquad \dot m = 0,
\label{eq:xaug}
\end{equation}
holding it constant across the horizon so that a window of measurements
constrains a single value rather than a free trajectory. Process noise
$\bm{w} \in \mathbb{R}^{13}$ is admitted on the physical states only;
adding noise to the mass channel would allow it to drift within the
window and would dissolve precisely the constraint that makes the
estimate meaningful.

Over a horizon of $N$ stages the estimator solves
\begin{align}
\min_{x_0,\, \bm{w}} \quad
& \sum_{k=0}^{N-1} \Bigl( \bigl\| \bm{y}_k - h(x_k) \bigr\|^2_{R}
+ \bigl\| \bm{w}_k \bigr\|^2_{Q} \Bigr)
+ \bigl\| x_0 - \bar{x}_0 \bigr\|^2_{Q_0} \nonumber \\
\mathrm{s.t.} \quad
& x_{k+1} = f\bigl( x_k,\, \bm{w}_k;\, p_{\mathrm{MHE}} \bigr), \nonumber \\
& \underline{m} \le m \le \overline{m},
\label{eq:mhe}
\end{align}
where $f$ is the discretisation of \eqref{eq:kin}--\eqref{eq:rot},
$h$ selects the measured states, and the arrival cost
$\| x_0 - \bar{x}_0 \|^2_{Q_0}$ anchors the start of the window to the
estimate propagated from the previous solve.

The payload geometry does \emph{not} appear in the decision vector. It is
supplied as a runtime parameter, alongside the measured input:
\begin{equation}
p_{\mathrm{MHE}} = \bigl[\, T,\; \bm{\tau},\; \dJ,\; c_x,\; c_y \,\bigr] \in \mathbb{R}^{7}.
\label{eq:pmhe}
\end{equation}
Within the estimator, therefore, $\dJ$ and $\cvec$ have exactly the same
status as the measured thrust and torque: known inputs, not unknowns.

We use $N = 20$ stages at $dt = \SI{0.1}{\second}$, a \SI{2.0}{\second}
horizon, and bound the mass to
$[\SI{1.961}{\kg},\, \SI{5.0}{\kg}]$. The lower bound is the physical
floor $0.95\,\mb$: the vehicle cannot weigh less than its own airframe.
That this bound is a \emph{hard constraint on a decision variable} is a
concrete advantage of the moving-horizon formulation over a recursive
filter such as an extended Kalman filter, which can be biased away from
non-physical estimates but cannot be forbidden from producing them.

\subsection{Ground-truth-free inversion of the CoM offset}\label{sec:model:conline}

Where rotor speeds are available, the centre-of-mass offset need not come
from the prior at all. In steady hover the eccentricity moment
\eqref{eq:tcom} is balanced by the inner loop, so
$\dot{\bm{\omega}} \approx 0$ and $\bm{\omega} \approx \bm{0}$ reduce
\eqref{eq:rot} to $\bm{\tau} = -\tcom$. Substituting \eqref{eq:tcom} and
solving componentwise,
\begin{equation}
c_x = -\,\frac{\tau_{\mathrm{pitch}}^{\mathrm{phys}}}{\Tphys},
\qquad
c_y = +\,\frac{\tau_{\mathrm{roll}}^{\mathrm{phys}}}{\Tphys}.
\label{eq:conline}
\end{equation}

This is an algebraic inversion --- one division --- not an optimisation,
and it introduces no decision variable. Three properties make it
admissible within an adaptive loop.

\begin{enumerate}
\item \textbf{It reads no model state.} Both numerator and denominator
come from \eqref{eq:phys}. The inversion therefore cannot close an
estimate-to-model feedback path, which is the failure mode analysed in
Section~\ref{sec:model:whynot}.

\item \textbf{It is immune to thrust calibration.} Substituting
$F_i = k_f \omega_i^2$ into \eqref{eq:phys} and \eqref{eq:conline}, the
rotor coefficient cancels,
\begin{equation}
c_x = -\frac{-\sum_i x_i k_f \omega_i^2}{\sum_i k_f \omega_i^2}
    = \frac{\sum_i x_i \omega_i^2}{\sum_i \omega_i^2},
\label{eq:kfcancel}
\end{equation}
leaving a quantity that depends only on rotor geometry and the
\emph{relative} distribution of rotor speeds. Any method that consumes an
absolute thrust inherits the calibration error of $k_f$; this one does
not.

\item \textbf{Channels are assigned by identifiability.} The offset
$\cvec$ possesses an independent observable and is estimated online,
whereas $r_z$ is annihilated by the cross product in \eqref{eq:tcom} and
cannot be identified from the torque channel at all. It must remain a
prior. The split is dictated by the physics, not chosen for convenience.
\end{enumerate}

Two filtering stages precede use of \eqref{eq:conline}, and they address
different error mechanisms. A \emph{gate} rejects samples taken outside
steady flight ($\|\bm{\omega}\| < \SI{0.15}{\radian\per\second}$,
$\|\bm{v}_{xy}\| < \SI{0.20}{\meter\per\second}$), because during
manoeuvres the measured torque contains angular-acceleration terms and is
not a trim torque at all --- a \emph{systematic} error that no amount of
averaging would remove. An exponential moving average with time constant
\SI{2}{\second} then suppresses \emph{random} error, which is
substantial: the numerator of \eqref{eq:conline} is a differential-mode
combination of rotor forces, in which the common-mode signal cancels and
the noise does not. Low-pass filtering costs no information here because
$\cvec$ is quasi-static once the load is grasped --- its signal lies at
DC, while the rotor noise does not. The mass, by contrast, changes as a
step at grasp and release, and could not be treated this way; that is
precisely why $m$ is handled by the optimisation \eqref{eq:mhe} and
$\cvec$ by a first-order filter.

\section{Control: How the Parameters Re-centre the Cost}\label{sec:model:ctrl}

The nonlinear model predictive controller carries the parameters as
runtime values,
\begin{equation}
p_{\mathrm{NMPC}} = \bigl[\, x_{\mathrm{ref}}(13),\; m,\; \dJ,\; \cvec(2),\; \bm{d}(3) \,\bigr] \in \mathbb{R}^{20},
\label{eq:pnmpc}
\end{equation}
where $\bm{d}$ is a lumped translational disturbance used only by the
$\mathcal{L}_1$ comparison baseline and identically zero under the
proposed method. Over a prediction horizon of $N$ stages it solves
\begin{align}
\min_{u_{0:N-1}} \quad
& \sum_{k=0}^{N-1} \Bigl( \bigl\| e_{\mathrm{track}}(x_k, x_{\mathrm{ref},k}) \bigr\|^2_{Q}
+ \bigl\| u_k - u_{\mathrm{hover}} \bigr\|^2_{R} \Bigr) \nonumber \\
& \qquad + \bigl\| e_{\mathrm{track}}(x_N, x_{\mathrm{ref},N}) \bigr\|^2_{P} \nonumber \\
\mathrm{s.t.} \quad
& x_{k+1} = f(x_k, u_k;\, p_{\mathrm{NMPC}}), \nonumber \\
& \underline{T} \le T_k \le \overline{T}, \quad
  |\tau_{k,i}| \le \overline{\tau}_i ,
\label{eq:nmpc}
\end{align}
applying only $u_0$ and re-solving at the next step.

\subsection{Why the input baseline must track the parameters}\label{sec:model:hover}

The second term of \eqref{eq:nmpc} penalises the deviation of the input
from a \emph{baseline} $u_{\mathrm{hover}}$, and that baseline is
recomputed from the live parameters at every solve rather than fixed when
the solver is generated:
\begin{equation}
T_h = m\,g,
\qquad
u_{\mathrm{hover}} = \bigl[\, T_h,\; c_y T_h,\; -c_x T_h,\; 0 \,\bigr]^{\!\top}.
\label{eq:hover}
\end{equation}

Consider first why a baseline is needed at all. Penalising $\|u\|^2_R$
outright would treat the thrust required merely to remain airborne as an
expense to be minimised, biasing the solution toward insufficient thrust.
The baseline recentres the penalty on the input that sustains
equilibrium, so that $R$ penalises only genuine control effort.

Consider next what happens when the baseline is \emph{stale}. Suppose
$u_{\mathrm{hover}}$ is computed from the unloaded mass while the vehicle
carries a parcel. The cost then pulls the thrust persistently toward a
value below what equilibrium requires, and the optimum settles where the
marginal tracking cost balances the marginal baseline-deviation cost ---
at a nonzero steady-state error. The behaviour has a counter-intuitive
signature that identifies it in practice: \emph{increasing $R$ makes the
error worse}, because $R$ is precisely the strength with which the
incorrect baseline pulls. No amount of weight tuning removes a bias of
this kind; only correcting the baseline does.

The torque entries of \eqref{eq:hover} are exactly $-\tcom$ from
\eqref{eq:tcom}, and they exist for the same reason. Hovering with an
eccentric load requires a constant trim torque
(Section~\ref{sec:model:tcom}); a cost that does not admit this pulls the
torque toward zero and produces the same class of steady-state bias in
the attitude channel. The two equations are complementary and neither
suffices alone: \eqref{eq:tcom} teaches the \emph{model} that the moment
exists, while \eqref{eq:hover} teaches the \emph{objective} that it is
necessary.

\subsection{Inner-loop bandwidth}\label{sec:model:gain}

Model consistency in the outer loop cannot restore inner-loop bandwidth.
The rate controller running on the flight controller is tuned for the
bare airframe; a payload that multiplies the roll and pitch inertia
lowers its effective bandwidth, and for heavy eccentric loads the
resulting attitude oscillation grows until the torque saturates. We
therefore rescale the rate gains at attachment by the inertia ratio,
\begin{equation}
K \leftarrow K_{\mathrm{nom}} \cdot
\min\!\Bigl( \frac{J_{xx} + \dJ}{J_{xx}},\; 5.0 \Bigr).
\label{eq:gain}
\end{equation}

The cap is not incidental. From \eqref{eq:numbers}, a \SI{0.3}{\kg}
payload already produces a ratio of $5.18$, so that payload sits at the
ceiling of what inner-loop retuning alone can absorb on this airframe.
Together with the torque budget \eqref{eq:trimtorque}, this defines a
feasibility boundary on payload mass and eccentricity that is
characterised experimentally in a later chapter.

\section{Why the Inertia and CoM Are Not Augmented into the State}\label{sec:model:whynot}

Augmenting $\dJ$ and $\cvec$ into the estimator's decision vector
\eqref{eq:xaug} --- yielding a seventeen-dimensional augmented state and
four online parameters instead of one --- is the obvious alternative, and
it is the design most readers will expect. We argue against it on four
grounds.

\paragraph{Asymmetric observability} The mass is excited directly by the
thrust--acceleration relation \eqref{eq:trans} and is observable in
ordinary near-hover flight, as argued in
Section~\ref{sec:model:trans}. The inertia increment acts only through
$\dot{\bm{\omega}}$ in \eqref{eq:rot}, and a delivery task --- spent
largely in near-hover, low-rate flight --- excites that channel weakly.
Admitting $\dJ$ as a free variable therefore permits the optimiser to
explain measurement noise using a nearly unobservable quantity, which
degrades rather than improves the estimate.

\paragraph{Competition for a single residual} Deriving the geometry from
the \emph{instantaneous} mass estimate closes a feedback path,
\begin{equation*}
m_{\mathrm{est}} \;\longrightarrow\; \dJ,\, \cvec
\;\longrightarrow\; \text{attitude model}
\;\longrightarrow\; m_{\mathrm{est}},
\end{equation*}
in which two quantities compete to explain the same residual. This
bootstrap was observed to fail structurally on light payloads: before
lift-off the mass estimate reads low \emph{legitimately}, since the load
still rests on the ground; that value is consumed as ``no geometry
present''; the geometry channel closes; and after lift-off the model is
mismatched and the estimate remains pinned at its lower bound. The remedy
was to sever the path --- letting $\mpay$ read a fixed operational prior
and never the running estimate --- not to add further free parameters.

\paragraph{Economy of parameterisation} Equations \eqref{eq:cxy} and
\eqref{eq:dj} obtain three numbers from one scalar. The two degrees of
freedom thereby saved are not discarded information: they are recovered
from a constraint, rigid attachment, that holds exactly.

\paragraph{The remaining prior may be weak} The single scalar $\mpay$
need not be accurate. With the geometry prior deliberately set
\SI{33}{\percent} wrong, the mass estimate still converges to within
\SI{0.3}{\percent}, so a nominal parcel mass --- a quantity a delivery
operator knows anyway --- suffices, and no ground-truth measurement of
the attachment is required. Where rotor speeds are available,
\eqref{eq:conline} removes the prior from the $\cvec$ path entirely,
leaving it to influence only $\dJ$.

\section{Declared Modelling Approximations}\label{sec:model:approx}

Three points are stated explicitly here rather than left to be
discovered.

\paragraph{Anisotropy of the inertia increment} The single scalar
\eqref{eq:dj} replaces the pair $(\dJ_{xx}, \dJ_{yy})$ of
\eqref{eq:djaxes} by their arithmetic mean. The resulting error is
$\tfrac{1}{2}\mu\,|r_x^2 - r_y^2|$, which vanishes exactly when
$r_x = r_y$ and remains below \SI{1}{\percent} of $J_{xx} + \dJ$ at the
offsets used in this work.

\paragraph{The yaw increment is neglected} The third expression in
\eqref{eq:djaxes}, $\dJ_{zz} = \mu(r_x^2 + r_y^2)$, is not added to
$J_{zz}$ in \eqref{eq:Jt}. At \SI{0.3}{\kg} and
$r_{xy} = \SI{0.109}{\meter}$ it amounts to
\SI{0.0031}{\kg\meter\squared}, or \SI{14.8}{\percent} of $J_{zz}$. This
is a known and non-negligible approximation, and it is retained only
because the yaw channel neither carries payload compensation nor
participates in the inversion \eqref{eq:conline}, so it does not affect
the results reported later. Removing it would require substituting
$J_{zz} + \mu(r_x^2 + r_y^2)$ into \eqref{eq:Jt} and would introduce no
new estimated quantity --- the correction is available at zero cost in
identification, and is omitted only for consistency with the
implementation as evaluated.

\paragraph{The vertical CoM offset is absent by construction} As shown in
Section~\ref{sec:model:tcom}, $c_z$ drops out of \eqref{eq:tcom} as an
exact algebraic consequence of collinearity, not as a truncation. It is
listed among the approximations only to forestall the opposite reading.
.
%%
%% REQUIRES in the preamble:
%%   amsmath, amssymb, booktabs, siunitx, bm
%%   (optional) \numberwithin{equation}{chapter}  for 3.1, 3.2, ...
%%
%% NOTE: the conference version of this material is sec_model.tex
%% (two-column IEEEtran, condensed). Keep both; they diverge on
%% purpose -- the thesis version carries the full derivations.
%% =====================================================================

% ---- Notation macros (move to the preamble if preferred) ------------
\providecommand{\dJ}{\Delta J}
\providecommand{\Jt}{J_t}
\providecommand{\Jb}{J_b}
\providecommand{\cvec}{\bm{c}}
\providecommand{\rp}{\bm{r}_p}
\providecommand{\rc}{\bm{r}_c}
\providecommand{\mpay}{\hat{m}_p}
\providecommand{\mb}{m_b}
\providecommand{\mt}{m_t}
\providecommand{\Tphys}{T_{\mathrm{phys}}}
\providecommand{\tauphys}{\bm{\tau}_{\mathrm{phys}}}
\providecommand{\tcom}{\bm{\tau}_{\mathrm{com}}}
\providecommand{\ev}{\bm{e}_3}

\chapter{Modeling and Parameter Structure}\label{ch:model}

A quadrotor that picks up a parcel does not merely become heavier. It
also becomes harder to rotate, and its centre of mass no longer lies on
the thrust axis. These are three distinct changes to the equations of
motion, and a controller that compensates only the first will fail in
ways that mass adaptation cannot explain.

This chapter derives the model that the estimator and the controller
share, and then makes explicit which of its parameters are identified
online and which are not. The central structural claim --- and the reason
the method needs so little sensing --- is that the three payload-induced
quantities are not three independent unknowns. Only the mass is a
decision variable of the estimator; the inertia increment and the
centre-of-mass offset are generated from it and from known contact
geometry by rigid-body algebra, at the cost of no additional degrees of
freedom.

Section~\ref{sec:model:frames} fixes frames and conventions,
Sections~\ref{sec:model:kin}--\ref{sec:model:dyn} derive the kinematics
and dynamics, and Section~\ref{sec:model:payload} derives the payload
parameterization. Sections~\ref{sec:model:phys}--\ref{sec:model:ctrl}
describe how the parameters are measured and consumed, and
Sections~\ref{sec:model:whynot}--\ref{sec:model:approx} defend the design
and declare its approximations.

\section{Frames, Notation, and Conventions}\label{sec:model:frames}

Two reference frames are used throughout. The world frame $\mathcal{W}$
is East--North--Up (ENU), so that gravity acts along $-z$. The body frame
$\mathcal{B}$ is Front--Left--Up (FLU) with its origin at the geometric
centre of the rotor disc --- \emph{not} at the centre of mass. This
choice is deliberate and consequential: the thrust vector passes through
the body origin by construction, so any offset between that origin and
the true centre of mass appears explicitly as a moment
(Section~\ref{sec:model:dyn}) rather than being absorbed silently into
the frame definition.

The state vector and the control input are
\begin{equation}
x = \bigl[\bm{p};\, \bm{v};\, \bm{q};\, \bm{\omega}\bigr] \in \mathbb{R}^{13},
\qquad
u = \bigl[T;\, \bm{\tau}\bigr] \in \mathbb{R}^{4},
\label{eq:state}
\end{equation}
where $\bm{p}$ and $\bm{v}$ are the position and velocity in
$\mathcal{W}$, $\bm{\omega}$ is the body rate in $\mathcal{B}$, $T$ is the
collective thrust along the body $z$ axis, and $\bm{\tau}$ is the body
torque. We abbreviate $\ev = [0,0,1]^{\!\top}$.

\subsection{Attitude representation}\label{sec:model:quat}

Attitude is represented by a unit quaternion under the \emph{Hamilton}
convention with the scalar part first,
\begin{equation}
\bm{q} = [q_w,\, q_x,\, q_y,\, q_z]^{\!\top},
\qquad \|\bm{q}\| = 1,
\label{eq:quatdef}
\end{equation}
encoding the rotation from $\mathcal{B}$ to $\mathcal{W}$. The
corresponding rotation matrix is
\begin{equation}
R(\bm{q}) =
\begin{bmatrix}
q_w^2 + q_x^2 - q_y^2 - q_z^2 & 2(q_xq_y - q_wq_z) & 2(q_xq_z + q_wq_y) \\[2pt]
2(q_xq_y + q_wq_z) & q_w^2 - q_x^2 + q_y^2 - q_z^2 & 2(q_yq_z - q_wq_x) \\[2pt]
2(q_xq_z - q_wq_y) & 2(q_yq_z + q_wq_x) & q_w^2 - q_x^2 - q_y^2 + q_z^2
\end{bmatrix}.
\label{eq:Rq}
\end{equation}

Three conventions have just been fixed, and all three matter. Hamilton
versus JPL determines the sign of the vector part in the quaternion
product; scalar-first versus scalar-last determines the layout of
\eqref{eq:Rq}; and resolving $\bm{\omega}$ in $\mathcal{B}$ rather than
$\mathcal{W}$ determines on which \emph{side} the product appears in the
kinematics \eqref{eq:quat} below. A mismatch in any one of them
integrates the attitude in the wrong direction while remaining
dimensionally consistent, which makes such errors unusually difficult to
detect. We state them here so that every later equation is unambiguous.

\section{Rigid-Body Kinematics}\label{sec:model:kin}

Translational kinematics is the definition of velocity,
\begin{equation}
\dot{\bm{p}} = \bm{v}.
\label{eq:kin}
\end{equation}

Attitude kinematics follows from the quaternion product of the attitude
with the pure quaternion formed from the body rate. Writing
$\bm{q} = (q_w, \bm{q}_v)$ with $\bm{q}_v = [q_x, q_y, q_z]^{\!\top}$,
and using the Hamilton product
$(a_w, \bm{a}_v) \otimes (b_w, \bm{b}_v) =
(a_w b_w - \bm{a}_v \!\cdot\! \bm{b}_v,\;
a_w \bm{b}_v + b_w \bm{a}_v + \bm{a}_v \!\times\! \bm{b}_v)$,
we obtain
\begin{equation}
\dot{\bm{q}} = \tfrac{1}{2}\, \bm{q} \otimes [0,\, \bm{\omega}]
= \tfrac{1}{2}
\begin{bmatrix}
-\bm{q}_v \cdot \bm{\omega} \\[2pt]
q_w \bm{\omega} + \bm{q}_v \times \bm{\omega}
\end{bmatrix}
= \tfrac{1}{2}\, \Xi(\bm{q})\, \bm{\omega},
\label{eq:quat}
\end{equation}
where the second equality expands the product and the third collects the
result into the $4 \times 3$ matrix
\begin{equation}
\Xi(\bm{q}) =
\begin{bmatrix}
-q_x & -q_y & -q_z \\
 q_w & -q_z &  q_y \\
 q_z &  q_w & -q_x \\
-q_y &  q_x &  q_w
\end{bmatrix}.
\label{eq:Xi}
\end{equation}
The two forms are algebraically identical; the implementation uses
\eqref{eq:Xi} because it is a plain matrix--vector product. Note that the
body rate enters on the \emph{right} of the product in \eqref{eq:quat}.
Had $\bm{\omega}$ been resolved in $\mathcal{W}$, the correct expression
would be $\dot{\bm{q}} = \tfrac{1}{2} [0, \bm{\omega}] \otimes \bm{q}$
instead, differing by a conjugation by $\bm{q}$.

Neither \eqref{eq:kin} nor \eqref{eq:quat} contains a mass, an inertia, or
any other physical parameter: kinematics is pure geometry. \emph{The
payload therefore enters the model only through the dynamics}, which
confines the entire parameter discussion of this chapter to the three
equations of Section~\ref{sec:model:dyn}.

\section{Rigid-Body Dynamics}\label{sec:model:dyn}

\subsection{Translational dynamics}\label{sec:model:trans}

Newton's second law in $\mathcal{W}$, with thrust rotated from the body
frame and a linear aerodynamic drag term, gives
\begin{equation}
\dot{\bm{v}} = \frac{1}{m}\Bigl( R(\bm{q})\, \ev\, T - k_d\, \bm{v} \Bigr) - g\, \ev,
\label{eq:trans}
\end{equation}
where $m$ is the \emph{composite} mass of airframe plus payload, $k_d$ is
a fixed drag coefficient, and $g$ is the gravitational acceleration.

Equation \eqref{eq:trans} is where the mass --- and \emph{only} the mass
--- enters, and it enters in exactly one role: as the scalar gain $1/m$
from thrust to acceleration. This single fact underwrites the
observability argument of the whole method. If a thrust value can be
obtained independently of the model, then the measured acceleration
determines $m$ uniquely by \eqref{eq:trans}, and it does so in ordinary
near-hover flight, with no dedicated excitation manoeuvre. The
qualification ``independently of the model'' is essential and is the
subject of Section~\ref{sec:model:phys}.

\subsection{Rotational dynamics}\label{sec:model:rot}

Let $\Jb = \mathrm{diag}(J_{xx}, J_{yy}, J_{zz})$ denote the
bare-airframe inertia, and let the payload raise the roll and pitch axes
by an increment $\dJ$. We write the resulting \emph{composite inertia} as
\begin{equation}
\Jt = \mathrm{diag}\bigl(J_{xx} + \dJ,\; J_{yy} + \dJ,\; J_{zz}\bigr),
\label{eq:Jt}
\end{equation}
the subscript matching the composite mass $\mt$ introduced in
Section~\ref{sec:model:payload}. Euler's equation for a rigid body,
written about the composite centre of mass and in the body frame, is then
\begin{equation}
\dot{\bm{\omega}} = \Jt^{-1}\Bigl( \bm{\tau} + \tcom(\cvec) - \bm{\omega} \times \Jt \bm{\omega} \Bigr).
\label{eq:rot}
\end{equation}
The term $\bm{\omega} \times \Jt \bm{\omega}$ is the usual gyroscopic
coupling. The term $\tcom$ is derived next.

That \eqref{eq:rot} is written \emph{about the composite centre of mass}
is not a stylistic choice: Euler's equation takes this form only when the
inertia tensor and the applied moments refer to the same point, and that
point is the centre of mass. Both $\Jt$ in Section~\ref{sec:model:payload}
and $\tcom$ below must therefore be computed about the composite CoM, not
about the body origin. Mixing the two reference points is a common and
silent modelling error.

\subsection{The thrust eccentricity moment}\label{sec:model:tcom}

Thrust acts at the rotor-disc centre, which by the frame convention of
Section~\ref{sec:model:frames} is the body origin. When a payload is
attached the composite centre of mass moves to some
$\rc = [c_x, c_y, c_z]^{\!\top} \neq \bm{0}$, and the thrust therefore no
longer passes through it. The moment about the CoM is the cross product
of the position of the application point \emph{relative to the CoM},
namely $-\rc$, with the force $\ev T$:
\begin{align}
\tcom(\cvec) &= (-\rc) \times \ev T
= -\begin{bmatrix} c_x \\ c_y \\ c_z \end{bmatrix}
\times \begin{bmatrix} 0 \\ 0 \\ T \end{bmatrix} \nonumber \\[4pt]
&= -\begin{bmatrix} c_y T - 0 \\ 0 - c_x T \\ 0 - 0 \end{bmatrix}
= \begin{bmatrix} -c_y T \\ +c_x T \\ 0 \end{bmatrix}.
\label{eq:tcom}
\end{align}

Two properties of \eqref{eq:tcom} deserve emphasis.

First, \textbf{$c_z$ cannot appear.} The third component of $\rc$
multiplies only terms that vanish, because a vertical CoM offset is
collinear with the body-$z$ thrust and the cross product of collinear
vectors is identically zero. Gravity, likewise, acts \emph{at} the centre
of mass and exerts no moment about it. The centre-of-mass offset is
therefore a two-parameter quantity \emph{by construction, not by
truncation} --- a distinction worth stating explicitly, since ``the
offset is three-dimensional, yet only two components are modelled'' is an
entirely natural objection to raise.

Second, \textbf{$\tcom$ is a constant moment in level hover.} It does not
average to zero, and it is not small: for the payload of
Section~\ref{sec:model:payload} it consumes roughly two thirds of the
available roll authority. Hovering on the spot with an off-centre load
therefore \emph{requires} a persistent trim torque. A model that omits
\eqref{eq:tcom} does not merely lose accuracy; it presents the controller
with a standing disturbance that has no explanation, which the controller
will then try to eliminate by manoeuvring. Section~\ref{sec:model:ctrl}
shows that the cost function must admit the same fact independently.

\subsection{Structure of the parameter entry points}\label{sec:model:entry}

Collecting \eqref{eq:trans}, \eqref{eq:rot} and \eqref{eq:tcom}, the model
has the property on which the rest of this thesis is built:

\begin{quote}
\emph{Each payload-dependent parameter has exactly one point of entry.}
The mass appears only as $1/m$ in the translational dynamics; the inertia
increment appears only inside $\Jt$ in the rotational dynamics; the
centre-of-mass offset appears only in the eccentricity moment.
\end{quote}

This separation is what makes the parameters individually attributable to
distinct physical channels --- the mass to the thrust--acceleration
relation, the inertia to angular acceleration, the offset to the
steady-state torque balance --- and it is exploited directly in
Sections~\ref{sec:model:mhe} and \ref{sec:model:whynot}.

\section{Payload-Induced Parameters}\label{sec:model:payload}

We now derive $\dJ$ and $\cvec$. The payload is modelled as a point mass
$\mpay$ rigidly attached at a known body-frame offset
$\rp = [r_x, r_y, r_z]^{\!\top}$ with $r_z < 0$ (the load hangs below the
airframe). Define the composite mass and the \emph{reduced mass}
\begin{equation}
\mt = \mb + \mpay,
\qquad
\mu = \frac{\mb\, \mpay}{\mb + \mpay} = \frac{\mb\, \mpay}{\mt},
\label{eq:mu}
\end{equation}
where $\mb$ is the bare-airframe mass.

\subsection{Centre-of-mass offset}\label{sec:model:com}

The composite centre of mass is the mass-weighted mean of the two
constituent positions. Taking the airframe centre of mass at the body
origin and the payload at $\rp$,
\begin{equation}
\rc = \frac{\mb \cdot \bm{0} + \mpay \cdot \rp}{\mt}
= \frac{\mpay}{\mt}\, \rp ,
\label{eq:rc}
\end{equation}
whose horizontal components are the pair required by \eqref{eq:tcom}:
\begin{equation}
\cvec = [c_x,\, c_y] = \frac{\mpay}{\mt}\, [r_x,\, r_y].
\label{eq:cxy}
\end{equation}

\subsection{Inertia increment and the reduced mass}\label{sec:model:paxis}

The inertia increment requires more care than the parallel-axis theorem
applied naively, because \emph{the centre of mass itself moves when the
payload is attached}. The theorem must be applied about the \emph{new}
centre of mass \eqref{eq:rc}, and both bodies contribute, since the
airframe no longer sits at the CoM either.

Consider a single axis and let $d$ denote the corresponding lever arm.
From \eqref{eq:rc}, the distances of the two bodies from the composite
centre of mass are
\begin{equation}
\underbrace{\bigl\| \bm{0} - \rc \bigr\| = \frac{\mpay}{\mt}\, d}_{\text{airframe}},
\qquad
\underbrace{\bigl\| \rp - \rc \bigr\| = \Bigl( 1 - \frac{\mpay}{\mt} \Bigr) d = \frac{\mb}{\mt}\, d}_{\text{payload}} .
\label{eq:dists}
\end{equation}
Note that the payload's distance is \emph{smaller} than $d$. Adding the
two parallel-axis contributions and simplifying,
\begin{align}
\Delta J^{(d)}
&= \mb \Bigl( \frac{\mpay}{\mt} d \Bigr)^{\!2}
 + \mpay \Bigl( \frac{\mb}{\mt} d \Bigr)^{\!2} \nonumber \\[4pt]
&= \frac{d^2}{\mt^2}\Bigl( \mb\, \mpay^2 + \mpay\, \mb^2 \Bigr)
 = \frac{d^2}{\mt^2}\, \mb\, \mpay \bigl( \mpay + \mb \bigr) \nonumber \\[4pt]
&= \frac{\mb\, \mpay}{\mt}\, d^2
 = \mu\, d^2 .
\label{eq:reduced}
\end{align}
The factor $(\mb + \mpay)$ produced by the sum cancels one power of $\mt$
in the denominator, and the two contributions collapse onto the reduced
mass \eqref{eq:mu} --- the same quantity that governs the two-body
problem in classical mechanics, and for the same reason: the inertia of a
rigid pair depends only on their \emph{relative} configuration.

Equation \eqref{eq:reduced} also identifies a natural error. Writing
$\Delta J^{(d)} = \mpay d^2$ --- treating the payload as orbiting a fixed
body origin, and forgetting both the airframe's own contribution and the
displacement of the CoM --- overestimates the increment by the factor
$\mt / \mb$, which is \SI{14.5}{\percent} at the payload considered
below. Two limiting cases confirm \eqref{eq:reduced}: as
$\mpay \!\ll\! \mb$ we recover $\mu \to \mpay$ (a light load barely moves
the CoM, and the naive expression becomes correct), while as
$\mpay \!\gg\! \mb$ we obtain $\mu \to \mb$ (a heavy load anchors the CoM
and it is the airframe that revolves about it).

Applying \eqref{eq:reduced} per axis with the appropriate arms gives
\begin{equation}
\dJ_{xx} = \mu\bigl(r_y^2 + r_z^2\bigr), \quad
\dJ_{yy} = \mu\bigl(r_x^2 + r_z^2\bigr), \quad
\dJ_{zz} = \mu\bigl(r_x^2 + r_y^2\bigr),
\label{eq:djaxes}
\end{equation}
and the model carries the roll/pitch mean as a single scalar,
\begin{equation}
\dJ = \tfrac{1}{2}\bigl(\dJ_{xx} + \dJ_{yy}\bigr)
    = \mu\Bigl( r_z^2 + \tfrac{1}{2}\bigl(r_x^2 + r_y^2\bigr) \Bigr).
\label{eq:dj}
\end{equation}
The consequences of using one scalar in place of the pair, and of
omitting $\dJ_{zz}$ from \eqref{eq:Jt}, are stated in
Section~\ref{sec:model:approx}.

\subsection{One unknown, three numbers}\label{sec:model:onescalar}

Equations \eqref{eq:cxy} and \eqref{eq:dj} produce all three
payload-dependent quantities --- $\dJ$, $c_x$, $c_y$ --- from a
\emph{single} unknown scalar $\mpay$, together with geometry $\rp$ that
is known from the gripper contact configuration. The rigid-attachment
constraint thus supplies two degrees of freedom that a naive augmentation
would otherwise have to identify from data. The estimator is not asked to
rediscover the parallel-axis theorem from measurement noise; it is given
the theorem and asked for one number.

This is the structural claim of the chapter, and
Table~\ref{tab:params} records the resulting parameter identities.
Section~\ref{sec:model:whynot} defends the choice against the obvious
alternative.

\begin{table}[t]
\caption{Model parameters, classified by identity. Only the mass is a
decision variable of the estimator; the inertia increment and the CoM
offset enter the estimator as \emph{known} runtime parameters, on the same
footing as the measured thrust and torque.}
\label{tab:params}
\centering
\footnotesize
\begin{tabular}{@{}l p{0.23\textwidth} c p{0.235\textwidth} p{0.17\textwidth}@{}}
\toprule
Symbol & Meaning & Dim. & Identity & Enters through \\
\midrule
$m$ & composite mass & 1 & \textbf{MHE decision variable} & \eqref{eq:trans}, \eqref{eq:hover} \\
$\dJ$ & roll/pitch inertia increment & 1 & derived, \eqref{eq:dj} & \eqref{eq:Jt}, \eqref{eq:gain} \\
$\cvec$ & CoM horizontal offset & 2 & derived \eqref{eq:cxy}, or inverted online \eqref{eq:conline} & \eqref{eq:tcom}, \eqref{eq:hover} \\
\midrule
$\mpay$ & payload mass (\emph{sole unknown}) & 1 & weak operational prior & \eqref{eq:cxy}, \eqref{eq:dj} \\
$\rp$ & payload offset in $\mathcal{B}$ & 3 & known (contact geometry) & \eqref{eq:cxy}, \eqref{eq:dj} \\
\midrule
$\mb$ & bare-airframe mass & 1 & calibration & \SI{2.0643}{\kg} \\
$\Jb$ & bare-airframe inertia & 3 & calibration & $\mathrm{diag}(0.0142,$ $0.0142, 0.0210)$ \\
$k_d$, $g$ & drag coefficient, gravity & 2 & constants & $0.05$, $9.81$ \\
\bottomrule
\end{tabular}
\end{table}

\subsection{Magnitudes on the experimental platform}\label{sec:model:magnitudes}

To fix ideas, consider a \SI{0.3}{\kg} payload attached at
$\rp = [0,\, 0.109,\, -0.47]\,$\si{\meter} on the platform of this work
($\mb = \SI{2.0643}{\kg}$, $J_{xx} = \SI{0.0142}{\kg\meter\squared}$,
roll/pitch torque authority \SI{0.5}{\newton\meter}). Equations
\eqref{eq:mu}, \eqref{eq:cxy} and \eqref{eq:dj} give
\begin{align}
\mu &= \SI{0.262}{\kg}, &
\dJ &= \SI{0.0594}{\kg\meter\squared}, \nonumber \\
\frac{J_{xx} + \dJ}{J_{xx}} &= 5.18, &
c_y &= \SI{13.8}{\milli\meter},
\label{eq:numbers}
\end{align}
and hence a constant hover trim torque
\begin{equation}
|c_y|\, \mt\, g = \SI{0.32}{\newton\meter} \approx \SI{64}{\percent}\ \text{of the roll authority.}
\label{eq:trimtorque}
\end{equation}

These numbers state the problem this thesis addresses more sharply than
any qualitative argument. A \SI{0.3}{\kg} parcel on a \SI{2.06}{\kg}
airframe is a \SI{15}{\percent} change of mass --- but simultaneously a
\emph{five-fold} change of roll inertia and a standing demand on two
thirds of the torque margin. Adapting the mass alone accounts for neither
of the latter two. Note also that the $r_z^2$ term contributes
\SI{97.4}{\percent} of $\dJ$ in \eqref{eq:dj}: the increment is dominated
by how far \emph{below} the airframe the load hangs, not by how far it is
off-centre.

\section{Model-Independent Reconstruction of the Physical Input}\label{sec:model:phys}

The thrust and torque appearing in \eqref{eq:trans} and \eqref{eq:rot}
must be known to the estimator. They are reconstructed from measured
rotor speeds $\omega_i$ and the known rotor positions $(x_i, y_i)$ via the
quadratic rotor model $F_i = k_f \omega_i^2$:
\begin{equation}
\Tphys = \sum_i F_i, \qquad
\tau_{\mathrm{roll}} = \sum_i y_i F_i, \qquad
\tau_{\mathrm{pitch}} = -\sum_i x_i F_i .
\label{eq:phys}
\end{equation}

It is worth being precise about why the \emph{commanded} thrust cannot be
used instead, since it is available for free and appears equivalent. Let
the controller compute its thrust command from an assumed mass
$\tilde{m}$, so that in hover $T_{\mathrm{cmd}} = \tilde{m} g$. Feeding
$T_{\mathrm{cmd}}$ into \eqref{eq:trans} and solving for the mass yields
$m = T_{\mathrm{cmd}} / g = \tilde{m}$ --- for \emph{any} $\tilde{m}$
whatsoever. The relation is satisfied identically, the residual carries no
information, and the mass is unobservable. The estimate would simply
reproduce whatever the controller already believed.

Equation \eqref{eq:phys} breaks this circularity because rotor speeds are
a physical measurement into which no model state enters. This is a
recurring principle in what follows: an adaptive loop must be fed by a
signal that is not itself a product of the adaptation.

\section{Estimation: The Mass as the Single Online Degree of Freedom}\label{sec:model:mhe}

\subsection{Augmented state and estimator formulation}\label{sec:model:mheform}

The moving-horizon estimator augments the mass into the state,
\begin{equation}
x_{\mathrm{aug}} = [x;\, m] \in \mathbb{R}^{14},
\qquad \dot m = 0,
\label{eq:xaug}
\end{equation}
holding it constant across the horizon so that a window of measurements
constrains a single value rather than a free trajectory. Process noise
$\bm{w} \in \mathbb{R}^{13}$ is admitted on the physical states only;
adding noise to the mass channel would allow it to drift within the
window and would dissolve precisely the constraint that makes the
estimate meaningful.

Over a horizon of $N$ stages the estimator solves
\begin{align}
\min_{x_0,\, \bm{w}} \quad
& \sum_{k=0}^{N-1} \Bigl( \bigl\| \bm{y}_k - h(x_k) \bigr\|^2_{R}
+ \bigl\| \bm{w}_k \bigr\|^2_{Q} \Bigr)
+ \bigl\| x_0 - \bar{x}_0 \bigr\|^2_{Q_0} \nonumber \\
\mathrm{s.t.} \quad
& x_{k+1} = f\bigl( x_k,\, \bm{w}_k;\, p_{\mathrm{MHE}} \bigr), \nonumber \\
& \underline{m} \le m \le \overline{m},
\label{eq:mhe}
\end{align}
where $f$ is the discretisation of \eqref{eq:kin}--\eqref{eq:rot},
$h$ selects the measured states, and the arrival cost
$\| x_0 - \bar{x}_0 \|^2_{Q_0}$ anchors the start of the window to the
estimate propagated from the previous solve.

The payload geometry does \emph{not} appear in the decision vector. It is
supplied as a runtime parameter, alongside the measured input:
\begin{equation}
p_{\mathrm{MHE}} = \bigl[\, T,\; \bm{\tau},\; \dJ,\; c_x,\; c_y \,\bigr] \in \mathbb{R}^{7}.
\label{eq:pmhe}
\end{equation}
Within the estimator, therefore, $\dJ$ and $\cvec$ have exactly the same
status as the measured thrust and torque: known inputs, not unknowns.

We use $N = 20$ stages at $dt = \SI{0.1}{\second}$, a \SI{2.0}{\second}
horizon, and bound the mass to
$[\SI{1.961}{\kg},\, \SI{5.0}{\kg}]$. The lower bound is the physical
floor $0.95\,\mb$: the vehicle cannot weigh less than its own airframe.
That this bound is a \emph{hard constraint on a decision variable} is a
concrete advantage of the moving-horizon formulation over a recursive
filter such as an extended Kalman filter, which can be biased away from
non-physical estimates but cannot be forbidden from producing them.

\subsection{Ground-truth-free inversion of the CoM offset}\label{sec:model:conline}

Where rotor speeds are available, the centre-of-mass offset need not come
from the prior at all. In steady hover the eccentricity moment
\eqref{eq:tcom} is balanced by the inner loop, so
$\dot{\bm{\omega}} \approx 0$ and $\bm{\omega} \approx \bm{0}$ reduce
\eqref{eq:rot} to $\bm{\tau} = -\tcom$. Substituting \eqref{eq:tcom} and
solving componentwise,
\begin{equation}
c_x = -\,\frac{\tau_{\mathrm{pitch}}^{\mathrm{phys}}}{\Tphys},
\qquad
c_y = +\,\frac{\tau_{\mathrm{roll}}^{\mathrm{phys}}}{\Tphys}.
\label{eq:conline}
\end{equation}

This is an algebraic inversion --- one division --- not an optimisation,
and it introduces no decision variable. Three properties make it
admissible within an adaptive loop.

\begin{enumerate}
\item \textbf{It reads no model state.} Both numerator and denominator
come from \eqref{eq:phys}. The inversion therefore cannot close an
estimate-to-model feedback path, which is the failure mode analysed in
Section~\ref{sec:model:whynot}.

\item \textbf{It is immune to thrust calibration.} Substituting
$F_i = k_f \omega_i^2$ into \eqref{eq:phys} and \eqref{eq:conline}, the
rotor coefficient cancels,
\begin{equation}
c_x = -\frac{-\sum_i x_i k_f \omega_i^2}{\sum_i k_f \omega_i^2}
    = \frac{\sum_i x_i \omega_i^2}{\sum_i \omega_i^2},
\label{eq:kfcancel}
\end{equation}
leaving a quantity that depends only on rotor geometry and the
\emph{relative} distribution of rotor speeds. Any method that consumes an
absolute thrust inherits the calibration error of $k_f$; this one does
not.

\item \textbf{Channels are assigned by identifiability.} The offset
$\cvec$ possesses an independent observable and is estimated online,
whereas $r_z$ is annihilated by the cross product in \eqref{eq:tcom} and
cannot be identified from the torque channel at all. It must remain a
prior. The split is dictated by the physics, not chosen for convenience.
\end{enumerate}

Two filtering stages precede use of \eqref{eq:conline}, and they address
different error mechanisms. A \emph{gate} rejects samples taken outside
steady flight ($\|\bm{\omega}\| < \SI{0.15}{\radian\per\second}$,
$\|\bm{v}_{xy}\| < \SI{0.20}{\meter\per\second}$), because during
manoeuvres the measured torque contains angular-acceleration terms and is
not a trim torque at all --- a \emph{systematic} error that no amount of
averaging would remove. An exponential moving average with time constant
\SI{2}{\second} then suppresses \emph{random} error, which is
substantial: the numerator of \eqref{eq:conline} is a differential-mode
combination of rotor forces, in which the common-mode signal cancels and
the noise does not. Low-pass filtering costs no information here because
$\cvec$ is quasi-static once the load is grasped --- its signal lies at
DC, while the rotor noise does not. The mass, by contrast, changes as a
step at grasp and release, and could not be treated this way; that is
precisely why $m$ is handled by the optimisation \eqref{eq:mhe} and
$\cvec$ by a first-order filter.

\section{Control: How the Parameters Re-centre the Cost}\label{sec:model:ctrl}

The nonlinear model predictive controller carries the parameters as
runtime values,
\begin{equation}
p_{\mathrm{NMPC}} = \bigl[\, x_{\mathrm{ref}}(13),\; m,\; \dJ,\; \cvec(2),\; \bm{d}(3) \,\bigr] \in \mathbb{R}^{20},
\label{eq:pnmpc}
\end{equation}
where $\bm{d}$ is a lumped translational disturbance used only by the
$\mathcal{L}_1$ comparison baseline and identically zero under the
proposed method. Over a prediction horizon of $N$ stages it solves
\begin{align}
\min_{u_{0:N-1}} \quad
& \sum_{k=0}^{N-1} \Bigl( \bigl\| e_{\mathrm{track}}(x_k, x_{\mathrm{ref},k}) \bigr\|^2_{Q}
+ \bigl\| u_k - u_{\mathrm{hover}} \bigr\|^2_{R} \Bigr) \nonumber \\
& \qquad + \bigl\| e_{\mathrm{track}}(x_N, x_{\mathrm{ref},N}) \bigr\|^2_{P} \nonumber \\
\mathrm{s.t.} \quad
& x_{k+1} = f(x_k, u_k;\, p_{\mathrm{NMPC}}), \nonumber \\
& \underline{T} \le T_k \le \overline{T}, \quad
  |\tau_{k,i}| \le \overline{\tau}_i ,
\label{eq:nmpc}
\end{align}
applying only $u_0$ and re-solving at the next step.

\subsection{Why the input baseline must track the parameters}\label{sec:model:hover}

The second term of \eqref{eq:nmpc} penalises the deviation of the input
from a \emph{baseline} $u_{\mathrm{hover}}$, and that baseline is
recomputed from the live parameters at every solve rather than fixed when
the solver is generated:
\begin{equation}
T_h = m\,g,
\qquad
u_{\mathrm{hover}} = \bigl[\, T_h,\; c_y T_h,\; -c_x T_h,\; 0 \,\bigr]^{\!\top}.
\label{eq:hover}
\end{equation}

Consider first why a baseline is needed at all. Penalising $\|u\|^2_R$
outright would treat the thrust required merely to remain airborne as an
expense to be minimised, biasing the solution toward insufficient thrust.
The baseline recentres the penalty on the input that sustains
equilibrium, so that $R$ penalises only genuine control effort.

Consider next what happens when the baseline is \emph{stale}. Suppose
$u_{\mathrm{hover}}$ is computed from the unloaded mass while the vehicle
carries a parcel. The cost then pulls the thrust persistently toward a
value below what equilibrium requires, and the optimum settles where the
marginal tracking cost balances the marginal baseline-deviation cost ---
at a nonzero steady-state error. The behaviour has a counter-intuitive
signature that identifies it in practice: \emph{increasing $R$ makes the
error worse}, because $R$ is precisely the strength with which the
incorrect baseline pulls. No amount of weight tuning removes a bias of
this kind; only correcting the baseline does.

The torque entries of \eqref{eq:hover} are exactly $-\tcom$ from
\eqref{eq:tcom}, and they exist for the same reason. Hovering with an
eccentric load requires a constant trim torque
(Section~\ref{sec:model:tcom}); a cost that does not admit this pulls the
torque toward zero and produces the same class of steady-state bias in
the attitude channel. The two equations are complementary and neither
suffices alone: \eqref{eq:tcom} teaches the \emph{model} that the moment
exists, while \eqref{eq:hover} teaches the \emph{objective} that it is
necessary.

\subsection{Inner-loop bandwidth}\label{sec:model:gain}

Model consistency in the outer loop cannot restore inner-loop bandwidth.
The rate controller running on the flight controller is tuned for the
bare airframe; a payload that multiplies the roll and pitch inertia
lowers its effective bandwidth, and for heavy eccentric loads the
resulting attitude oscillation grows until the torque saturates. We
therefore rescale the rate gains at attachment by the inertia ratio,
\begin{equation}
K \leftarrow K_{\mathrm{nom}} \cdot
\min\!\Bigl( \frac{J_{xx} + \dJ}{J_{xx}},\; 5.0 \Bigr).
\label{eq:gain}
\end{equation}

The cap is not incidental. From \eqref{eq:numbers}, a \SI{0.3}{\kg}
payload already produces a ratio of $5.18$, so that payload sits at the
ceiling of what inner-loop retuning alone can absorb on this airframe.
Together with the torque budget \eqref{eq:trimtorque}, this defines a
feasibility boundary on payload mass and eccentricity that is
characterised experimentally in a later chapter.

\section{Why the Inertia and CoM Are Not Augmented into the State}\label{sec:model:whynot}

Augmenting $\dJ$ and $\cvec$ into the estimator's decision vector
\eqref{eq:xaug} --- yielding a seventeen-dimensional augmented state and
four online parameters instead of one --- is the obvious alternative, and
it is the design most readers will expect. We argue against it on four
grounds.

\paragraph{Asymmetric observability} The mass is excited directly by the
thrust--acceleration relation \eqref{eq:trans} and is observable in
ordinary near-hover flight, as argued in
Section~\ref{sec:model:trans}. The inertia increment acts only through
$\dot{\bm{\omega}}$ in \eqref{eq:rot}, and a delivery task --- spent
largely in near-hover, low-rate flight --- excites that channel weakly.
Admitting $\dJ$ as a free variable therefore permits the optimiser to
explain measurement noise using a nearly unobservable quantity, which
degrades rather than improves the estimate.

\paragraph{Competition for a single residual} Deriving the geometry from
the \emph{instantaneous} mass estimate closes a feedback path,
\begin{equation*}
m_{\mathrm{est}} \;\longrightarrow\; \dJ,\, \cvec
\;\longrightarrow\; \text{attitude model}
\;\longrightarrow\; m_{\mathrm{est}},
\end{equation*}
in which two quantities compete to explain the same residual. This
bootstrap was observed to fail structurally on light payloads: before
lift-off the mass estimate reads low \emph{legitimately}, since the load
still rests on the ground; that value is consumed as ``no geometry
present''; the geometry channel closes; and after lift-off the model is
mismatched and the estimate remains pinned at its lower bound. The remedy
was to sever the path --- letting $\mpay$ read a fixed operational prior
and never the running estimate --- not to add further free parameters.

\paragraph{Economy of parameterisation} Equations \eqref{eq:cxy} and
\eqref{eq:dj} obtain three numbers from one scalar. The two degrees of
freedom thereby saved are not discarded information: they are recovered
from a constraint, rigid attachment, that holds exactly.

\paragraph{The remaining prior may be weak} The single scalar $\mpay$
need not be accurate. With the geometry prior deliberately set
\SI{33}{\percent} wrong, the mass estimate still converges to within
\SI{0.3}{\percent}, so a nominal parcel mass --- a quantity a delivery
operator knows anyway --- suffices, and no ground-truth measurement of
the attachment is required. Where rotor speeds are available,
\eqref{eq:conline} removes the prior from the $\cvec$ path entirely,
leaving it to influence only $\dJ$.

\section{Declared Modelling Approximations}\label{sec:model:approx}

Three points are stated explicitly here rather than left to be
discovered.

\paragraph{Anisotropy of the inertia increment} The single scalar
\eqref{eq:dj} replaces the pair $(\dJ_{xx}, \dJ_{yy})$ of
\eqref{eq:djaxes} by their arithmetic mean. The resulting error is
$\tfrac{1}{2}\mu\,|r_x^2 - r_y^2|$, which vanishes exactly when
$r_x = r_y$ and remains below \SI{1}{\percent} of $J_{xx} + \dJ$ at the
offsets used in this work.

\paragraph{The yaw increment is neglected} The third expression in
\eqref{eq:djaxes}, $\dJ_{zz} = \mu(r_x^2 + r_y^2)$, is not added to
$J_{zz}$ in \eqref{eq:Jt}. At \SI{0.3}{\kg} and
$r_{xy} = \SI{0.109}{\meter}$ it amounts to
\SI{0.0031}{\kg\meter\squared}, or \SI{14.8}{\percent} of $J_{zz}$. This
is a known and non-negligible approximation, and it is retained only
because the yaw channel neither carries payload compensation nor
participates in the inversion \eqref{eq:conline}, so it does not affect
the results reported later. Removing it would require substituting
$J_{zz} + \mu(r_x^2 + r_y^2)$ into \eqref{eq:Jt} and would introduce no
new estimated quantity --- the correction is available at zero cost in
identification, and is omitted only for consistency with the
implementation as evaluated.

\paragraph{The vertical CoM offset is absent by construction} As shown in
Section~\ref{sec:model:tcom}, $c_z$ drops out of \eqref{eq:tcom} as an
exact algebraic consequence of collinearity, not as a truncation. It is
listed among the approximations only to forestall the opposite reading.
.
%%
%% REQUIRES in the preamble:
%%   amsmath, amssymb, booktabs, siunitx, bm
%%   (optional) \numberwithin{equation}{chapter}  for 3.1, 3.2, ...
%%
%% NOTE: the conference version of this material is sec_model.tex
%% (two-column IEEEtran, condensed). Keep both; they diverge on
%% purpose -- the thesis version carries the full derivations.
%% =====================================================================

% ---- Notation macros (move to the preamble if preferred) ------------
\providecommand{\dJ}{\Delta J}
\providecommand{\Jt}{J_t}
\providecommand{\Jb}{J_b}
\providecommand{\cvec}{\bm{c}}
\providecommand{\rp}{\bm{r}_p}
\providecommand{\rc}{\bm{r}_c}
\providecommand{\mpay}{\hat{m}_p}
\providecommand{\mb}{m_b}
\providecommand{\mt}{m_t}
\providecommand{\Tphys}{T_{\mathrm{phys}}}
\providecommand{\tauphys}{\bm{\tau}_{\mathrm{phys}}}
\providecommand{\tcom}{\bm{\tau}_{\mathrm{com}}}
\providecommand{\ev}{\bm{e}_3}

\chapter{Modeling and Parameter Structure}\label{ch:model}
\begin{spacing}{1.5}
\setlength{\parskip}{0.22in}

A quadrotor that picks up a parcel does not merely become heavier. It
also becomes harder to rotate, and its centre of mass no longer lies on
the thrust axis. These are three distinct changes to the equations of
motion, and a controller that compensates only the first will fail in
ways that mass adaptation cannot explain.

This chapter derives the model that the estimator and the controller
share, and then makes explicit which of its parameters are identified
online and which are not. The central structural claim --- and the reason
the method needs so little sensing --- is that the three payload-induced
quantities are not three independent unknowns. Only the composite mass is
a decision variable of the estimator. The parcel mass is predicted from
the increment above the unloaded vehicle model; the inertia increment and
centre-of-mass offset are then generated from this prediction and contact
geometry by rigid-body algebra, at the cost of no additional optimisation
degrees of freedom. No task-specific parcel-mass value is admitted to the
MHE; the controller-side safety scales retained in revision
\texttt{f2264b3} are separated explicitly later in this chapter.

Section~\ref{sec:model:frames} fixes frames and conventions,
Sections~\ref{sec:model:kin}--\ref{sec:model:dyn} derive the kinematics
and dynamics, and Section~\ref{sec:model:payload} derives the payload
parameterization. Sections~\ref{sec:model:phys}--\ref{sec:model:ctrl}
describe how the parameters are measured and consumed, and
Sections~\ref{sec:model:whynot}--\ref{sec:model:approx} defend the design
and declare its approximations.

\section{Frames, Notation, and Conventions}\label{sec:model:frames}

Two reference frames are used throughout. The world frame $\mathcal{W}$
is East--North--Up (ENU), so that gravity acts along $-z$. The body frame
$\mathcal{B}$ is Front--Left--Up (FLU) with its origin at the geometric
centre of the rotor disc --- \emph{not} at the centre of mass. This
choice is deliberate and consequential: the thrust vector passes through
the body origin by construction, so any offset between that origin and
the true centre of mass appears explicitly as a moment
(Section~\ref{sec:model:dyn}) rather than being absorbed silently into
the frame definition.

The state vector and the control input are
\begin{equation}
x = \bigl[\bm{p};\, \bm{v};\, \bm{q};\, \bm{\omega}\bigr] \in \mathbb{R}^{13},
\qquad
u = \bigl[T;\, \bm{\tau}\bigr] \in \mathbb{R}^{4},
\label{eq:state}
\end{equation}
where $\bm{p}$ and $\bm{v}$ are the position and velocity in
$\mathcal{W}$, $\bm{\omega}$ is the body rate in $\mathcal{B}$, $T$ is the
collective thrust along the body $z$ axis, and $\bm{\tau}$ is the body
torque. We abbreviate $\ev = [0,0,1]^{\!\top}$.

\subsection{Attitude representation}\label{sec:model:quat}

Attitude is represented by a unit quaternion under the \emph{Hamilton}
convention with the scalar part first,
\begin{equation}
\bm{q} = [q_w,\, q_x,\, q_y,\, q_z]^{\!\top},
\qquad \|\bm{q}\| = 1,
\label{eq:quatdef}
\end{equation}
encoding the rotation from $\mathcal{B}$ to $\mathcal{W}$. The
corresponding rotation matrix is
\begin{equation}
R(\bm{q}) =
\begin{bmatrix}
q_w^2 + q_x^2 - q_y^2 - q_z^2 & 2(q_xq_y - q_wq_z) & 2(q_xq_z + q_wq_y) \\[2pt]
2(q_xq_y + q_wq_z) & q_w^2 - q_x^2 + q_y^2 - q_z^2 & 2(q_yq_z - q_wq_x) \\[2pt]
2(q_xq_z - q_wq_y) & 2(q_yq_z + q_wq_x) & q_w^2 - q_x^2 - q_y^2 + q_z^2
\end{bmatrix}.
\label{eq:Rq}
\end{equation}

Three conventions have just been fixed, and all three matter. Hamilton
versus JPL determines the sign of the vector part in the quaternion
product; scalar-first versus scalar-last determines the layout of
\eqref{eq:Rq}; and resolving $\bm{\omega}$ in $\mathcal{B}$ rather than
$\mathcal{W}$ determines on which \emph{side} the product appears in the
kinematics \eqref{eq:quat} below. A mismatch in any one of them
integrates the attitude in the wrong direction while remaining
dimensionally consistent, which makes such errors unusually difficult to
detect. We state them here so that every later equation is unambiguous.

\section{Rigid-Body Kinematics}\label{sec:model:kin}

Translational kinematics is the definition of velocity,
\begin{equation}
\dot{\bm{p}} = \bm{v}.
\label{eq:kin}
\end{equation}

Attitude kinematics follows from the quaternion product of the attitude
with the pure quaternion formed from the body rate. Writing
$\bm{q} = (q_w, \bm{q}_v)$ with $\bm{q}_v = [q_x, q_y, q_z]^{\!\top}$,
and using the Hamilton product
$(a_w, \bm{a}_v) \otimes (b_w, \bm{b}_v) =
(a_w b_w - \bm{a}_v \!\cdot\! \bm{b}_v,\;
a_w \bm{b}_v + b_w \bm{a}_v + \bm{a}_v \!\times\! \bm{b}_v)$,
we obtain
\begin{equation}
\dot{\bm{q}} = \tfrac{1}{2}\, \bm{q} \otimes [0,\, \bm{\omega}]
= \tfrac{1}{2}
\begin{bmatrix}
-\bm{q}_v \cdot \bm{\omega} \\[2pt]
q_w \bm{\omega} + \bm{q}_v \times \bm{\omega}
\end{bmatrix}
= \tfrac{1}{2}\, \Xi(\bm{q})\, \bm{\omega},
\label{eq:quat}
\end{equation}
where the second equality expands the product and the third collects the
result into the $4 \times 3$ matrix
\begin{equation}
\Xi(\bm{q}) =
\begin{bmatrix}
-q_x & -q_y & -q_z \\
 q_w & -q_z &  q_y \\
 q_z &  q_w & -q_x \\
-q_y &  q_x &  q_w
\end{bmatrix}.
\label{eq:Xi}
\end{equation}
The two forms are algebraically identical; the implementation uses
\eqref{eq:Xi} because it is a plain matrix--vector product. Note that the
body rate enters on the \emph{right} of the product in \eqref{eq:quat}.
Had $\bm{\omega}$ been resolved in $\mathcal{W}$, the correct expression
would be $\dot{\bm{q}} = \tfrac{1}{2} [0, \bm{\omega}] \otimes \bm{q}$
instead, differing by a conjugation by $\bm{q}$.

Neither \eqref{eq:kin} nor \eqref{eq:quat} contains a mass, an inertia, or
any other physical parameter: kinematics is pure geometry. \emph{The
payload therefore enters the model only through the dynamics}, which
confines the entire parameter discussion of this chapter to the three
equations of Section~\ref{sec:model:dyn}.

\section{Rigid-Body Dynamics}\label{sec:model:dyn}

\subsection{Translational dynamics}\label{sec:model:trans}

Newton's second law in $\mathcal{W}$, with thrust rotated from the body
frame and a linear aerodynamic drag term, gives
\begin{equation}
\dot{\bm{v}} = \frac{1}{m}\Bigl( R(\bm{q})\, \ev\, T - k_d\, \bm{v} \Bigr) - g\, \ev,
\label{eq:trans}
\end{equation}
where $m$ is the \emph{composite} mass of airframe plus payload, $k_d$ is
a fixed drag coefficient, and $g$ is the gravitational acceleration.

Equation \eqref{eq:trans} is where the mass --- and \emph{only} the mass
--- enters, and it enters in exactly one role: as the scalar gain $1/m$
from thrust to acceleration. This single fact underwrites the
observability argument of the whole method. If a thrust value can be
obtained independently of the model, then the measured acceleration
determines $m$ uniquely by \eqref{eq:trans}, and it does so in ordinary
near-hover flight, with no dedicated excitation manoeuvre. The
qualification ``independently of the model'' is essential and is the
subject of Section~\ref{sec:model:phys}.

\subsection{Rotational dynamics}\label{sec:model:rot}

Let $\Jb = \mathrm{diag}(J_{xx}, J_{yy}, J_{zz})$ denote the
bare-airframe inertia, and let the payload raise the roll and pitch axes
by an increment $\dJ$. We write the resulting \emph{composite inertia} as
\begin{equation}
\Jt = \mathrm{diag}\bigl(J_{xx} + \dJ,\; J_{yy} + \dJ,\; J_{zz}\bigr),
\label{eq:Jt}
\end{equation}
the subscript matching the composite mass $\mt$ introduced in
Section~\ref{sec:model:payload}. Euler's equation for a rigid body,
written about the composite centre of mass and in the body frame, is then
\begin{equation}
\dot{\bm{\omega}} = \Jt^{-1}\Bigl( \bm{\tau} + \tcom(\cvec) - \bm{\omega} \times \Jt \bm{\omega} \Bigr).
\label{eq:rot}
\end{equation}
The term $\bm{\omega} \times \Jt \bm{\omega}$ is the usual gyroscopic
coupling. The term $\tcom$ is derived next.

That \eqref{eq:rot} is written \emph{about the composite centre of mass}
is not a stylistic choice: Euler's equation takes this form only when the
inertia tensor and the applied moments refer to the same point, and that
point is the centre of mass. Both $\Jt$ in Section~\ref{sec:model:payload}
and $\tcom$ below must therefore be computed about the composite CoM, not
about the body origin. Mixing the two reference points is a common and
silent modelling error.

\subsection{The thrust eccentricity moment}\label{sec:model:tcom}

Thrust acts at the rotor-disc centre, which by the frame convention of
Section~\ref{sec:model:frames} is the body origin. When a payload is
attached the composite centre of mass moves to some
$\rc = [c_x, c_y, c_z]^{\!\top} \neq \bm{0}$, and the thrust therefore no
longer passes through it. The moment about the CoM is the cross product
of the position of the application point \emph{relative to the CoM},
namely $-\rc$, with the force $\ev T$:
\begin{align}
\tcom(\cvec) &= (-\rc) \times \ev T
= -\begin{bmatrix} c_x \\ c_y \\ c_z \end{bmatrix}
\times \begin{bmatrix} 0 \\ 0 \\ T \end{bmatrix} \nonumber \\[4pt]
&= -\begin{bmatrix} c_y T - 0 \\ 0 - c_x T \\ 0 - 0 \end{bmatrix}
= \begin{bmatrix} -c_y T \\ +c_x T \\ 0 \end{bmatrix}.
\label{eq:tcom}
\end{align}

Two properties of \eqref{eq:tcom} deserve emphasis.

First, \textbf{$c_z$ cannot appear.} The third component of $\rc$
multiplies only terms that vanish, because a vertical CoM offset is
collinear with the body-$z$ thrust and the cross product of collinear
vectors is identically zero. Gravity, likewise, acts \emph{at} the centre
of mass and exerts no moment about it. The centre-of-mass offset is
therefore a two-parameter quantity \emph{by construction, not by
truncation} --- a distinction worth stating explicitly, since ``the
offset is three-dimensional, yet only two components are modelled'' is an
entirely natural objection to raise.

Second, \textbf{$\tcom$ is a constant moment in level hover.} It does not
average to zero, and it is not small: for the payload of
Section~\ref{sec:model:payload} it consumes roughly two thirds of the
available roll authority. Hovering on the spot with an off-centre load
therefore \emph{requires} a persistent trim torque. A model that omits
\eqref{eq:tcom} does not merely lose accuracy; it presents the controller
with a standing disturbance that has no explanation, which the controller
will then try to eliminate by manoeuvring. Section~\ref{sec:model:ctrl}
shows that the cost function must admit the same fact independently.

\subsection{Structure of the parameter entry points}\label{sec:model:entry}

Collecting \eqref{eq:trans}, \eqref{eq:rot} and \eqref{eq:tcom}, the model
has the property on which the rest of this thesis is built:

\begin{quote}
\emph{Each payload-dependent parameter has exactly one point of entry.}
The mass appears only as $1/m$ in the translational dynamics; the inertia
increment appears only inside $\Jt$ in the rotational dynamics; the
centre-of-mass offset appears only in the eccentricity moment.
\end{quote}

This separation is what makes the parameters individually attributable to
distinct physical channels --- the mass to the thrust--acceleration
relation, the inertia to angular acceleration, the offset to the
steady-state torque balance --- and it is exploited directly in
Sections~\ref{sec:model:mhe} and \ref{sec:model:whynot}.

\section{Payload-Induced Parameters}\label{sec:model:payload}

We now derive $\dJ$ and $\cvec$. The payload is modelled as a point mass
$\mpay$ rigidly attached at a known body-frame offset
$\rp = [r_x, r_y, r_z]^{\!\top}$ with $r_z < 0$ (the load hangs below the
airframe). Define the composite mass and the \emph{reduced mass}
\begin{equation}
\mt = \mb + \mpay,
\qquad
\mu = \frac{\mb\, \mpay}{\mb + \mpay} = \frac{\mb\, \mpay}{\mt},
\label{eq:mu}
\end{equation}
where $\mb$ is the bare-airframe mass.

\subsection{Centre-of-mass offset}\label{sec:model:com}

The composite centre of mass is the mass-weighted mean of the two
constituent positions. Taking the airframe centre of mass at the body
origin and the payload at $\rp$,
\begin{equation}
\rc = \frac{\mb \cdot \bm{0} + \mpay \cdot \rp}{\mt}
= \frac{\mpay}{\mt}\, \rp ,
\label{eq:rc}
\end{equation}
whose horizontal components are the pair required by \eqref{eq:tcom}:
\begin{equation}
\cvec = [c_x,\, c_y] = \frac{\mpay}{\mt}\, [r_x,\, r_y].
\label{eq:cxy}
\end{equation}

\subsection{Inertia increment and the reduced mass}\label{sec:model:paxis}

The inertia increment requires more care than the parallel-axis theorem
applied naively, because \emph{the centre of mass itself moves when the
payload is attached}. The theorem must be applied about the \emph{new}
centre of mass \eqref{eq:rc}, and both bodies contribute, since the
airframe no longer sits at the CoM either.

Consider a single axis and let $d$ denote the corresponding lever arm.
From \eqref{eq:rc}, the distances of the two bodies from the composite
centre of mass are
\begin{equation}
\underbrace{\bigl\| \bm{0} - \rc \bigr\| = \frac{\mpay}{\mt}\, d}_{\text{airframe}},
\qquad
\underbrace{\bigl\| \rp - \rc \bigr\| = \Bigl( 1 - \frac{\mpay}{\mt} \Bigr) d = \frac{\mb}{\mt}\, d}_{\text{payload}} .
\label{eq:dists}
\end{equation}
Note that the payload's distance is \emph{smaller} than $d$. Adding the
two parallel-axis contributions and simplifying,
\begin{align}
\Delta J^{(d)}
&= \mb \Bigl( \frac{\mpay}{\mt} d \Bigr)^{\!2}
 + \mpay \Bigl( \frac{\mb}{\mt} d \Bigr)^{\!2} \nonumber \\[4pt]
&= \frac{d^2}{\mt^2}\Bigl( \mb\, \mpay^2 + \mpay\, \mb^2 \Bigr)
 = \frac{d^2}{\mt^2}\, \mb\, \mpay \bigl( \mpay + \mb \bigr) \nonumber \\[4pt]
&= \frac{\mb\, \mpay}{\mt}\, d^2
 = \mu\, d^2 .
\label{eq:reduced}
\end{align}
The factor $(\mb + \mpay)$ produced by the sum cancels one power of $\mt$
in the denominator, and the two contributions collapse onto the reduced
mass \eqref{eq:mu} --- the same quantity that governs the two-body
problem in classical mechanics, and for the same reason: the inertia of a
rigid pair depends only on their \emph{relative} configuration.

Equation \eqref{eq:reduced} also identifies a natural error. Writing
$\Delta J^{(d)} = \mpay d^2$ --- treating the payload as orbiting a fixed
body origin, and forgetting both the airframe's own contribution and the
displacement of the CoM --- overestimates the increment by the factor
$\mt / \mb$, which is \SI{14.5}{\percent} at the payload considered
below. Two limiting cases confirm \eqref{eq:reduced}: as
$\mpay \!\ll\! \mb$ we recover $\mu \to \mpay$ (a light load barely moves
the CoM, and the naive expression becomes correct), while as
$\mpay \!\gg\! \mb$ we obtain $\mu \to \mb$ (a heavy load anchors the CoM
and it is the airframe that revolves about it).

Applying \eqref{eq:reduced} per axis with the appropriate arms gives
\begin{equation}
\dJ_{xx} = \mu\bigl(r_y^2 + r_z^2\bigr), \quad
\dJ_{yy} = \mu\bigl(r_x^2 + r_z^2\bigr), \quad
\dJ_{zz} = \mu\bigl(r_x^2 + r_y^2\bigr),
\label{eq:djaxes}
\end{equation}
and the model carries the roll/pitch mean as a single scalar,
\begin{equation}
\dJ = \tfrac{1}{2}\bigl(\dJ_{xx} + \dJ_{yy}\bigr)
    = \mu\Bigl( r_z^2 + \tfrac{1}{2}\bigl(r_x^2 + r_y^2\bigr) \Bigr).
\label{eq:dj}
\end{equation}
The consequences of using one scalar in place of the pair, and of
omitting $\dJ_{zz}$ from \eqref{eq:Jt}, are stated in
Section~\ref{sec:model:approx}.

\subsection{One unknown, three numbers}\label{sec:model:onescalar}

Equations \eqref{eq:cxy} and \eqref{eq:dj} produce all three
payload-dependent quantities --- $\dJ$, $c_x$, $c_y$ --- from a
\emph{single} online scalar $\mpay$, together with geometry $\rp$ from
the gripper contact configuration. The scalar is not a supplied prior:
it is computed after each MHE solve as the non-negative increment of the
estimated composite mass above the unloaded airframe model. The
rigid-attachment constraint thus supplies two degrees of freedom that a
naive augmentation would otherwise have to identify from data. The
estimator is not asked to rediscover the parallel-axis theorem from
measurement noise; it is given the theorem and asked for one mass.

This is the structural claim of the chapter, and
Table~\ref{tab:params} records the resulting parameter identities.
Section~\ref{sec:model:whynot} defends the choice against the obvious
alternative.

\begin{table}[t]
\caption{Model parameters for the prior-free coupled baseline. Only
composite mass is an MHE decision variable; parcel mass and the coupled
geometry terms are deterministic functions or independent physical-input
inversions.}
\label{tab:params}
\centering
\footnotesize
\begin{tabular}{@{}l p{0.23\textwidth} c p{0.235\textwidth} p{0.17\textwidth}@{}}
\toprule
Symbol & Meaning & Dim. & Identity & Enters through \\
\midrule
$m$ & composite mass & 1 & \textbf{MHE decision variable} & \eqref{eq:trans}, \eqref{eq:hover} \\
$\dJ$ & roll/pitch inertia increment & 1 & derived online from $\hat m_p$ & \eqref{eq:Jt}, \eqref{eq:gain} \\
$\cvec$ & CoM horizontal offset & 2 & derived \eqref{eq:cxy}, or inverted online \eqref{eq:conline} & \eqref{eq:tcom}, \eqref{eq:hover} \\
\midrule
$\mpay$ & predicted payload mass & 1 & $\max(\hat m-\mb,0)$; no parcel prior & \eqref{eq:cxy}, \eqref{eq:dj} \\
$\rp$ & payload offset in $\mathcal{B}$ & 3 & known (contact geometry) & \eqref{eq:cxy}, \eqref{eq:dj} \\
\midrule
$\mb$ & bare-airframe mass & 1 & calibration & \SI{2.0643}{\kg} \\
$\Jb$ & bare-airframe inertia & 3 & calibration & $\mathrm{diag}(0.0142,$ $0.0142, 0.0210)$ \\
$k_d$, $g$ & drag coefficient, gravity & 2 & constants & $0.05$, $9.81$ \\
\bottomrule
\end{tabular}
\end{table}

\subsection{Magnitudes on the experimental platform}\label{sec:model:magnitudes}

To fix ideas, consider a \SI{0.3}{\kg} payload attached at
$\rp = [0,\, 0.109,\, -0.47]\,$\si{\meter} on the platform of this work
($\mb = \SI{2.0643}{\kg}$, $J_{xx} = \SI{0.0142}{\kg\meter\squared}$,
roll/pitch torque authority \SI{0.5}{\newton\meter}). Equations
\eqref{eq:mu}, \eqref{eq:cxy} and \eqref{eq:dj} give
\begin{align}
\mu &= \SI{0.262}{\kg}, &
\dJ &= \SI{0.0594}{\kg\meter\squared}, \nonumber \\
\frac{J_{xx} + \dJ}{J_{xx}} &= 5.18, &
c_y &= \SI{13.8}{\milli\meter},
\label{eq:numbers}
\end{align}
and hence a constant hover trim torque
\begin{equation}
|c_y|\, \mt\, g = \SI{0.32}{\newton\meter} \approx \SI{64}{\percent}\ \text{of the roll authority.}
\label{eq:trimtorque}
\end{equation}

These numbers state the problem this thesis addresses more sharply than
any qualitative argument. A \SI{0.3}{\kg} parcel on a \SI{2.06}{\kg}
airframe is a \SI{15}{\percent} change of mass --- but simultaneously a
\emph{five-fold} change of roll inertia and a standing demand on two
thirds of the torque margin. Adapting the mass alone accounts for neither
of the latter two. Note also that the $r_z^2$ term contributes
\SI{97.4}{\percent} of $\dJ$ in \eqref{eq:dj}: the increment is dominated
by how far \emph{below} the airframe the load hangs, not by how far it is
off-centre.

\section{Model-Independent Reconstruction of the Physical Input}\label{sec:model:phys}

The thrust and torque appearing in \eqref{eq:trans} and \eqref{eq:rot}
must be known to the estimator. They are reconstructed from measured
rotor speeds $\omega_i$ and the known rotor positions $(x_i, y_i)$ via the
quadratic rotor model $F_i = k_f \omega_i^2$:
\begin{equation}
\Tphys = \sum_i F_i, \qquad
\tau_{\mathrm{roll}} = \sum_i y_i F_i, \qquad
\tau_{\mathrm{pitch}} = -\sum_i x_i F_i .
\label{eq:phys}
\end{equation}

It is worth being precise about why the \emph{commanded} thrust cannot be
used instead, since it is available for free and appears equivalent. Let
the controller compute its thrust command from an assumed mass
$\tilde{m}$, so that in hover $T_{\mathrm{cmd}} = \tilde{m} g$. Feeding
$T_{\mathrm{cmd}}$ into \eqref{eq:trans} and solving for the mass yields
$m = T_{\mathrm{cmd}} / g = \tilde{m}$ --- for \emph{any} $\tilde{m}$
whatsoever. The relation is satisfied identically, the residual carries no
information, and the mass is unobservable. The estimate would simply
reproduce whatever the controller already believed.

Equation \eqref{eq:phys} breaks this circularity because rotor speeds are
a physical measurement into which no model state enters. This is a
recurring principle in what follows: an adaptive loop must be fed by a
signal that is not itself a product of the adaptation.

\section{Estimation: The Mass as the Single Online Degree of Freedom}\label{sec:model:mhe}

\subsection{Augmented state and estimator formulation}\label{sec:model:mheform}

The moving-horizon estimator augments the mass into the state,
\begin{equation}
x_{\mathrm{aug}} = [x;\, m] \in \mathbb{R}^{14},
\qquad \dot m = 0,
\label{eq:xaug}
\end{equation}
holding it constant across the horizon so that a window of measurements
constrains a single value rather than a free trajectory. Process noise
$\bm{w} \in \mathbb{R}^{13}$ is admitted on the physical states only;
adding noise to the mass channel would allow it to drift within the
window and would dissolve precisely the constraint that makes the
estimate meaningful.

Over a horizon of $N$ stages the estimator solves
\begin{align}
\min_{x_0,\, \bm{w}} \quad
& \sum_{k=0}^{N-1} \Bigl( \bigl\| \bm{y}_k - h(x_k) \bigr\|^2_{R}
+ \bigl\| \bm{w}_k \bigr\|^2_{Q} \Bigr)
+ \bigl\| x_0 - \bar{x}_0 \bigr\|^2_{Q_0} \nonumber \\
\mathrm{s.t.} \quad
& x_{k+1} = f\bigl( x_k,\, \bm{w}_k;\, p_{\mathrm{MHE}} \bigr), \nonumber \\
& \underline{m} \le m \le \overline{m},
\label{eq:mhe}
\end{align}
where $f$ is the discretisation of \eqref{eq:kin}--\eqref{eq:rot},
$h$ selects the measured states, and the arrival cost
$\| x_0 - \bar{x}_0 \|^2_{Q_0}$ anchors the start of the window to the
estimate propagated from the previous solve.

The payload geometry does \emph{not} appear in the decision vector. In
the coupled no-task-mass MHE profile, the contact vector is supplied as a
runtime parameter alongside the measured input:
\begin{equation}
p_{\mathrm{MHE}} = \bigl[\, T,\; \bm{\tau},\; r_x,\; r_y,\; r_z \,\bigr] \in \mathbb{R}^{7}.
\label{eq:pmhe}
\end{equation}
Within the estimator, $\dJ(m)$ and $\cvec(m)$ are evaluated from the
same mass decision used by the translational model. This preserves the
mass derivative through both translational and rotational dynamics
without adding a payload-mass prior or additional decision variable.

We use $N = 20$ stages at $dt = \SI{0.1}{\second}$, a \SI{2.0}{\second}
horizon, and bound the mass to
$[\SI{0.2}{\kg},\, \SI{5.0}{\kg}]$. The lower value contains no
airframe information; it exists only to keep the $1/m$ dynamics away from
the singularity. In particular, the former bound
$0.95\,\mb=\SI{1.961}{\kg}$ is not used, because it reveals the
unloaded answer to the optimiser. The mass entry of the arrival matrix is
$Q_{0,m}=10^{-6}$ rather than $0.1$. This retains a numerically non-zero
Hessian entry while corresponding to an uninformative standard deviation
of approximately \SI{1000}{\kg}.

\subsection{Prior-free initialisation and parcel-mass prediction}
\label{sec:model:noprior}

Before NMPC takeover, rotor-speed reconstruction supplies a physical
hover input. When the altitude and thrust gates indicate airborne
support, the first mass guess is
\begin{equation}
m^{(0)}=\operatorname{clip}\!\left(
\frac{\Tphys}{g},\;\underline m,\;\overline m\right).
\label{eq:massseed}
\end{equation}
If a reliable physical-thrust sample is unavailable, the fallback is the
midpoint of the numerical box, not $\mb$. Equation~\eqref{eq:massseed}
is an initial iterate rather than a prior: its arrival weight is
$10^{-6}$, and subsequent MHE solutions are determined by the complete
state and physical-input window. The pre-offboard estimate is withheld
from NMPC until 20 consecutive \SI{10}{Hz} values have a standard
deviation below \SI{0.02}{\kg}, with a \SI{30}{s} timeout that is
reported rather than silently accepted.

After pickup, the estimator continues to solve for total mass $\hat m_t$.
The algorithmic parcel prediction supplied to the payload map is
\begin{equation}
\hat m_p=\max\!\left(\hat m_t-\mb,0\right).
\label{eq:payloadprediction}
\end{equation}
Here $\mb$ is the calibrated mass of the reusable vehicle, like $\Jb$
or $k_f$; it contains no information about the parcel selected for the
mission. Neither \eqref{eq:massseed}, the arrival cost nor the mass box
uses a parcel-specific mass. At drop the estimate in
\eqref{eq:payloadprediction} is not assigned zero. The geometry mode is
removed and the same MHE is left to drive $\hat m_t$ continuously back
toward the unloaded solution, so return to zero payload mass is an
observable estimator outcome rather than a state-machine command.

\subsection{Ground-truth-free inversion of the CoM offset}\label{sec:model:conline}

Where rotor speeds are available, the centre-of-mass offset need not come
from the prior at all. In steady hover the eccentricity moment
\eqref{eq:tcom} is balanced by the inner loop, so
$\dot{\bm{\omega}} \approx 0$ and $\bm{\omega} \approx \bm{0}$ reduce
\eqref{eq:rot} to $\bm{\tau} = -\tcom$. Substituting \eqref{eq:tcom} and
solving componentwise,
\begin{equation}
c_x = -\,\frac{\tau_{\mathrm{pitch}}^{\mathrm{phys}}}{\Tphys},
\qquad
c_y = +\,\frac{\tau_{\mathrm{roll}}^{\mathrm{phys}}}{\Tphys}.
\label{eq:conline}
\end{equation}

This is an algebraic inversion --- one division --- not an optimisation,
and it introduces no decision variable. Three properties make it
admissible within an adaptive loop.

\begin{enumerate}
\item \textbf{It reads no model state.} Both numerator and denominator
come from \eqref{eq:phys}. The inversion therefore cannot close an
estimate-to-model feedback path, which is the failure mode analysed in
Section~\ref{sec:model:whynot}.

\item \textbf{It is immune to thrust calibration.} Substituting
$F_i = k_f \omega_i^2$ into \eqref{eq:phys} and \eqref{eq:conline}, the
rotor coefficient cancels,
\begin{equation}
c_x = -\frac{-\sum_i x_i k_f \omega_i^2}{\sum_i k_f \omega_i^2}
    = \frac{\sum_i x_i \omega_i^2}{\sum_i \omega_i^2},
\label{eq:kfcancel}
\end{equation}
leaving a quantity that depends only on rotor geometry and the
\emph{relative} distribution of rotor speeds. Any method that consumes an
absolute thrust inherits the calibration error of $k_f$; this one does
not.

\item \textbf{Channels are assigned by identifiability.} The offset
$\cvec$ possesses an independent observable and is estimated online,
whereas $r_z$ is annihilated by the cross product in \eqref{eq:tcom} and
cannot be identified from the torque channel at all. It must remain a
prior. The split is dictated by the physics, not chosen for convenience.
\end{enumerate}

Two filtering stages precede use of \eqref{eq:conline}, and they address
different error mechanisms. A \emph{gate} rejects samples taken outside
steady flight ($\|\bm{\omega}\| < \SI{0.15}{\radian\per\second}$,
$\|\bm{v}_{xy}\| < \SI{0.20}{\meter\per\second}$), because during
manoeuvres the measured torque contains angular-acceleration terms and is
not a trim torque at all --- a \emph{systematic} error that no amount of
averaging would remove. An exponential moving average with time constant
\SI{2}{\second} then suppresses \emph{random} error, which is
substantial: the numerator of \eqref{eq:conline} is a differential-mode
combination of rotor forces, in which the common-mode signal cancels and
the noise does not. Low-pass filtering costs no information here because
$\cvec$ is quasi-static once the load is grasped --- its signal lies at
DC, while the rotor noise does not. The mass, by contrast, changes as a
step at grasp and release, and could not be treated this way; that is
precisely why $m$ is handled by the optimisation \eqref{eq:mhe} and
$\cvec$ by a first-order filter.

\section{Control: How the Parameters Re-centre the Cost}\label{sec:model:ctrl}

The nonlinear model predictive controller carries the parameters as
runtime values,
\begin{equation}
p_{\mathrm{NMPC}} = \bigl[\, x_{\mathrm{ref}}(13),\; m,\; \dJ,\; \cvec(2),\; \bm{d}(3) \,\bigr] \in \mathbb{R}^{20},
\label{eq:pnmpc}
\end{equation}
where $\bm{d}$ is a lumped translational disturbance used only by the
$\mathcal{L}_1$ comparison baseline and identically zero under the
proposed method. Over a prediction horizon of $N$ stages it solves
\begin{align}
\min_{u_{0:N-1}} \quad
& \sum_{k=0}^{N-1} \Bigl( \bigl\| e_{\mathrm{track}}(x_k, x_{\mathrm{ref},k}) \bigr\|^2_{Q}
+ \bigl\| u_k - u_{\mathrm{hover}} \bigr\|^2_{R} \Bigr) \nonumber \\
& \qquad + \bigl\| e_{\mathrm{track}}(x_N, x_{\mathrm{ref},N}) \bigr\|^2_{P} \nonumber \\
\mathrm{s.t.} \quad
& x_{k+1} = f(x_k, u_k;\, p_{\mathrm{NMPC}}), \nonumber \\
& \underline{T} \le T_k \le \overline{T}, \quad
  |\tau_{k,i}| \le \overline{\tau}_i ,
\label{eq:nmpc}
\end{align}
applying only $u_0$ and re-solving at the next step.

\subsection{Why the input baseline must track the parameters}\label{sec:model:hover}

The second term of \eqref{eq:nmpc} penalises the deviation of the input
from a \emph{baseline} $u_{\mathrm{hover}}$, and that baseline is
recomputed from the live parameters at every solve rather than fixed when
the solver is generated:
\begin{equation}
T_h = m\,g,
\qquad
u_{\mathrm{hover}} = \bigl[\, T_h,\; c_y T_h,\; -c_x T_h,\; 0 \,\bigr]^{\!\top}.
\label{eq:hover}
\end{equation}

Consider first why a baseline is needed at all. Penalising $\|u\|^2_R$
outright would treat the thrust required merely to remain airborne as an
expense to be minimised, biasing the solution toward insufficient thrust.
The baseline recentres the penalty on the input that sustains
equilibrium, so that $R$ penalises only genuine control effort.

Consider next what happens when the baseline is \emph{stale}. Suppose
$u_{\mathrm{hover}}$ is computed from the unloaded mass while the vehicle
carries a parcel. The cost then pulls the thrust persistently toward a
value below what equilibrium requires, and the optimum settles where the
marginal tracking cost balances the marginal baseline-deviation cost ---
at a nonzero steady-state error. The behaviour has a counter-intuitive
signature that identifies it in practice: \emph{increasing $R$ makes the
error worse}, because $R$ is precisely the strength with which the
incorrect baseline pulls. No amount of weight tuning removes a bias of
this kind; only correcting the baseline does.

The torque entries of \eqref{eq:hover} are exactly $-\tcom$ from
\eqref{eq:tcom}, and they exist for the same reason. Hovering with an
eccentric load requires a constant trim torque
(Section~\ref{sec:model:tcom}); a cost that does not admit this pulls the
torque toward zero and produces the same class of steady-state bias in
the attitude channel. The two equations are complementary and neither
suffices alone: \eqref{eq:tcom} teaches the \emph{model} that the moment
exists, while \eqref{eq:hover} teaches the \emph{objective} that it is
necessary.

\subsection{Inner-loop bandwidth}\label{sec:model:gain}

Model consistency in the outer loop cannot restore inner-loop bandwidth.
The rate controller running on the flight controller is tuned for the
bare airframe; a payload that multiplies the roll and pitch inertia
lowers its effective bandwidth, and for heavy eccentric loads the
resulting attitude oscillation grows until the torque saturates. We
therefore rescale the rate gains at attachment by the inertia ratio,
\begin{equation}
K \leftarrow K_{\mathrm{nom}} \cdot
\min\!\Bigl( \frac{J_{xx} + \dJ}{J_{xx}},\; 5.0 \Bigr).
\label{eq:gain}
\end{equation}

The cap is not incidental. From \eqref{eq:numbers}, a \SI{0.3}{\kg}
payload already produces a ratio of $5.18$, so that payload sits at the
ceiling of what inner-loop retuning alone can absorb on this airframe.
Together with the torque budget \eqref{eq:trimtorque}, this defines a
feasibility boundary on payload mass and eccentricity that is
characterised experimentally in a later chapter.

\section{Why the Inertia and CoM Are Not Augmented into the State}\label{sec:model:whynot}

Augmenting $\dJ$ and $\cvec$ into the estimator's decision vector
\eqref{eq:xaug} --- yielding a seventeen-dimensional augmented state and
four online parameters instead of one --- is the obvious alternative, and
it is the design most readers will expect. We argue against it on four
grounds.

\paragraph{Asymmetric observability} The mass is excited directly by the
thrust--acceleration relation \eqref{eq:trans} and is observable in
ordinary near-hover flight, as argued in
Section~\ref{sec:model:trans}. The inertia increment acts only through
$\dot{\bm{\omega}}$ in \eqref{eq:rot}, and a delivery task --- spent
largely in near-hover, low-rate flight --- excites that channel weakly.
Admitting $\dJ$ as a free variable therefore permits the optimiser to
explain measurement noise using a nearly unobservable quantity, which
degrades rather than improves the estimate.

\paragraph{Competition for a single residual} Deriving geometry outside
the optimiser from a filtered or ratcheted mass estimate closes a feedback path,
\begin{equation*}
m_{\mathrm{est}} \;\longrightarrow\; \dJ,\, \cvec
\;\longrightarrow\; \text{attitude model}
\;\longrightarrow\; m_{\mathrm{est}},
\end{equation*}
in which two quantities compete to explain the same residual. This
bootstrap was observed to fail structurally on light payloads: before
lift-off the mass estimate reads low \emph{legitimately}, since the load
still rests on the ground; that value is consumed as ``no geometry
present''; the geometry channel closes; and after lift-off the model is
mismatched and the estimate remains pinned at its lower bound. The
coupled MHE option removes that outer fixed-point iteration: the measured
contact vector $\rp$ is supplied and $J(m)$ and $\cvec(m)$ are evaluated
inside the same symbolic dynamics used to differentiate the mass
decision. For the independent online CoM path, \eqref{eq:conline} uses
only physical rotor quantities. Revision \texttt{f2264b3} additionally
implements a separate NMPC-side engineering path in which the published
mass estimate drives a bounded $\dJ$ tracker. That path is deliberately
not presented as mathematically equivalent to the internally coupled
MHE: its ratchet and floor trade unbiased two-way tracking for transient
robustness and must be evaluated as controller safety conditioning.

\paragraph{Economy of parameterisation} Equations \eqref{eq:cxy} and
\eqref{eq:dj} obtain three numbers from one scalar. The two degrees of
freedom thereby saved are not discarded information: they are recovered
from a constraint, rigid attachment, that holds exactly.

\paragraph{Mass-information boundary} A platform load envelope may be
used to define actuator limits, startup gain scaling and safety
termination, but it is not inserted as $\mpay$ and does not seed, anchor
or lower-bound the mass optimisation. In the coupled MHE configuration,
$\rp$ occupies the three geometry parameter slots and the envelope is not
consumed by the estimator model. Equation~\eqref{eq:payloadprediction}
supplies the payload-mass value used by the coupled inertia map, while
\eqref{eq:conline} supplies a separately observable CoM check. The
deployed controller nevertheless retains mass-valued safety constants for
the attach-to-lift transient. Consequently, the precise claim is ``no
task-specific parcel mass enters the MHE'', not ``the complete flight
stack contains no payload-mass scale whatsoever''.

\section{Declared Modelling Approximations}\label{sec:model:approx}

Three points are stated explicitly here rather than left to be
discovered.

\paragraph{Anisotropy of the inertia increment} The single scalar
\eqref{eq:dj} replaces the pair $(\dJ_{xx}, \dJ_{yy})$ of
\eqref{eq:djaxes} by their arithmetic mean. The resulting error is
$\tfrac{1}{2}\mu\,|r_x^2 - r_y^2|$, which vanishes exactly when
$r_x = r_y$ and remains below \SI{1}{\percent} of $J_{xx} + \dJ$ at the
offsets used in this work.

\paragraph{The yaw increment is neglected} The third expression in
\eqref{eq:djaxes}, $\dJ_{zz} = \mu(r_x^2 + r_y^2)$, is not added to
$J_{zz}$ in \eqref{eq:Jt}. At \SI{0.3}{\kg} and
$r_{xy} = \SI{0.109}{\meter}$ it amounts to
\SI{0.0031}{\kg\meter\squared}, or \SI{14.8}{\percent} of $J_{zz}$. This
is a known and non-negligible approximation, and it is retained only
because the yaw channel neither carries payload compensation nor
participates in the inversion \eqref{eq:conline}, so it does not affect
the results reported later. Removing it would require substituting
$J_{zz} + \mu(r_x^2 + r_y^2)$ into \eqref{eq:Jt} and would introduce no
new estimated quantity --- the correction is available at zero cost in
identification, and is omitted only for consistency with the
implementation as evaluated.

\paragraph{The vertical CoM offset is absent by construction} As shown in
Section~\ref{sec:model:tcom}, $c_z$ drops out of \eqref{eq:tcom} as an
exact algebraic consequence of collinearity, not as a truncation. It is
listed among the approximations only to forestall the opposite reading.

\end{spacing}
\newpage
